\documentclass[11pt]{article}

\usepackage[margin=1in]{geometry}
\usepackage{amsmath,amssymb,amsthm}
\usepackage{enumitem}
\usepackage{xcolor}
\usepackage{hyperref}

\hypersetup{
  colorlinks=true,
  linkcolor=blue!60!black,
  urlcolor=blue!60!black
}

\newcommand{\R}{\mathbb{R}}
\newcommand{\N}{\mathbb{N}}
\newcommand{\B}{\mathcal{B}}
\newcommand{\F}{\mathcal{F}}
\newcommand{\Prob}{\mathbb{P}}
\newcommand{\E}{\mathbb{E}}
\newcommand{\1}{\mathbf{1}}
\newcommand{\Knowledge}[1]{%
  \par\smallskip\noindent
  \textbf{Knowledge used.}
  {\small\color{blue!60!black}\raggedright #1\par}
  \par\smallskip
}

\title{Chapter 2 Exercises}
\author{\textsc{Solutions to Rick Durrett's}\\
\textit{Probability: Theory and Examples}}
\date{\today}

\begin{document}
\maketitle

\noindent
These solutions use the local Chapter 2 LLM/Viki knowledge set in
\path{Probability/LLM_Viki_Dataset/Chapter2_Knowledge}.  The cited
knowledge IDs refer to entries in
\path{chapter2_knowledge_pieces.jsonl}.

\section*{Exercise 2.1.1}

\noindent\textbf{Question.}
Suppose $(X_1,\ldots,X_n)$ has density $f(x_1,\ldots,x_n)$, that is
\[
  \Prob((X_1,\ldots,X_n)\in A)=\int_A f(x)\,dx,\qquad A\in\B(\R^n).
\]
If $f(x)$ can be written as $g_1(x_1)\cdots g_n(x_n)$, where the
$g_m\ge 0$ are measurable, then $X_1,\ldots,X_n$ are independent.  Note
that the $g_m$ are not assumed to be probability densities.

\Knowledge{
\path{durrett_ch2_joint_law_product_measure};
\path{durrett_ch2_pi_lambda_independence};
\path{durrett_ch2_rv_sigma_independence}.
}

\noindent\textbf{Solution.}
Let
\[
  c_i=\int_{\R}g_i(x)\,dx.
\]
By Tonelli's theorem,
\[
  1=\int_{\R^n}f(x)\,dx
   =\prod_{i=1}^n \int_{\R}g_i(x_i)\,dx_i
   =\prod_{i=1}^n c_i.
\]
Thus each $c_i$ is finite and positive.  Define $h_i=g_i/c_i$.  Then
each $h_i$ is a probability density and
\[
  f(x_1,\ldots,x_n)=\prod_{i=1}^n h_i(x_i).
\]
For rectangles $A_1\times\cdots\times A_n$,
\[
\begin{aligned}
  \Prob(X_1\in A_1,\ldots,X_n\in A_n)
  &=\int_{A_1\times\cdots\times A_n}\prod_{i=1}^n h_i(x_i)\,dx  \\
  &=\prod_{i=1}^n \int_{A_i}h_i(x_i)\,dx_i
   =\prod_{i=1}^n \Prob(X_i\in A_i).
\end{aligned}
\]
The rectangle class is a pi-system generating $\B(\R^n)$, so the
factorization extends to the generated sigma-fields.  Hence
$X_1,\ldots,X_n$ are independent.

\section*{Exercise 2.1.2}

\noindent\textbf{Question.}
Suppose $X_1,\ldots,X_n$ are random variables that take values in
countable sets $S_1,\ldots,S_n$.  Then in order for
$X_1,\ldots,X_n$ to be independent, it is sufficient that whenever
$x_i\in S_i$,
\[
  \Prob(X_1=x_1,\ldots,X_n=x_n)=\prod_{i=1}^n \Prob(X_i=x_i).
\]

\Knowledge{
\path{durrett_ch2_event_independence};
\path{durrett_ch2_pi_lambda_independence};
\path{durrett_ch2_joint_law_product_measure}.
}

\noindent\textbf{Solution.}
Let $A_i\subseteq S_i$.  Since the $S_i$ are countable, Tonelli's theorem
for nonnegative sums gives
\[
\begin{aligned}
  \Prob(X_1\in A_1,\ldots,X_n\in A_n)
  &=\sum_{x_1\in A_1}\cdots\sum_{x_n\in A_n}
    \Prob(X_1=x_1,\ldots,X_n=x_n)\\
  &=\sum_{x_1\in A_1}\cdots\sum_{x_n\in A_n}
    \prod_{i=1}^n \Prob(X_i=x_i)\\
  &=\prod_{i=1}^n \sum_{x_i\in A_i}\Prob(X_i=x_i)
   =\prod_{i=1}^n \Prob(X_i\in A_i).
\end{aligned}
\]
Thus the defining product formula holds for arbitrary subsets of the
countable state spaces, and $X_1,\ldots,X_n$ are independent.

\section*{Exercise 2.1.3}

\noindent\textbf{Question.}
Let $\rho(x,y)$ be a metric.
\begin{enumerate}[label=(\roman*)]
\item Suppose $h$ is differentiable with $h(0)=0$, $h'(x)>0$ for
$x>0$, and $h'(x)$ decreasing on $[0,\infty)$.  Then $h(\rho(x,y))$ is a
metric.
\item $h(x)=x/(x+1)$ satisfies the hypotheses in (i).
\end{enumerate}

\Knowledge{
\path{durrett_ch2_construct_independent_variables}.
}

\noindent\textbf{Solution.}
Since $h'(x)>0$ for $x>0$ and $h(0)=0$, $h(t)\ge 0$ for $t\ge 0$, and
$h(t)=0$ if and only if $t=0$.  Therefore
$d(x,y)=h(\rho(x,y))$ is nonnegative and equals $0$ exactly when $x=y$.
Symmetry follows immediately from symmetry of $\rho$.

It remains to prove the triangle inequality.  Since $h'$ is decreasing,
$h$ is concave on $[0,\infty)$.  For $a,b\ge 0$,
\[
  h(a+b)-h(a)=\int_a^{a+b}h'(u)\,du
  \le \int_0^b h'(u)\,du
  =h(b).
\]
Thus $h(a+b)\le h(a)+h(b)$.  Since $\rho(x,z)\le \rho(x,y)+\rho(y,z)$
and $h$ is increasing,
\[
  h(\rho(x,z))
  \le h(\rho(x,y)+\rho(y,z))
  \le h(\rho(x,y))+h(\rho(y,z)).
\]
So $h\circ\rho$ is a metric.

For $h(x)=x/(x+1)$,
\[
  h(0)=0,\qquad h'(x)=\frac{1}{(x+1)^2}>0,\qquad
  h''(x)=-\frac{2}{(x+1)^3}<0.
\]
Hence $h'$ is decreasing and the hypotheses in (i) hold.

\section*{Exercise 2.1.4}

\noindent\textbf{Question.}
Let $\Omega=(0,1)$, $\F$ be the Borel sets, and $\Prob$ be Lebesgue
measure.  Let
\[
  X_n(\omega)=\sin(2\pi n\omega),\qquad n=1,2,\ldots.
\]
Show that $X_n$ are uncorrelated but not independent.

\Knowledge{
\path{durrett_ch2_expectation_factorization};
\path{durrett_ch2_rv_sigma_independence};
\path{durrett_ch2_event_independence}.
}

\noindent\textbf{Solution.}
For every $n$,
\[
  \E X_n=\int_0^1 \sin(2\pi n\omega)\,d\omega=0.
\]
If $m\ne n$, the trigonometric orthogonality identity gives
\[
  \E(X_mX_n)
  =\int_0^1 \sin(2\pi m\omega)\sin(2\pi n\omega)\,d\omega=0.
\]
Thus $\E(X_mX_n)=\E X_m\,\E X_n$ for $m\ne n$, so the variables are
uncorrelated.

They are not independent.  Indeed,
\[
  X_2^2=\sin^2(4\pi\omega)
  =4\sin^2(2\pi\omega)\bigl(1-\sin^2(2\pi\omega)\bigr)
  =4X_1^2(1-X_1^2).
\]
Thus $X_2^2$ is $\sigma(X_1)$-measurable.  If $X_1$ and $X_2$ were
independent, then $X_2^2$ would be independent of $\sigma(X_1)$ and
hence independent of itself.  A random variable independent of itself
must be almost surely constant, but $4X_1^2(1-X_1^2)$ is not almost
surely constant.  Therefore $X_1$ and $X_2$ are not independent, so the
sequence is not independent.

\section*{Exercise 2.1.5}

\noindent\textbf{Question.}
\begin{enumerate}[label=(\roman*)]
\item Show that if $X$ and $Y$ are independent with distributions $\mu$
and $\nu$, then
\[
  \Prob(X+Y=0)=\sum_y \mu(\{-y\})\nu(\{y\}).
\]
\item Conclude that if $X$ has continuous distribution, then
$\Prob(X=Y)=0$ for independent $X$ and $Y$.
\end{enumerate}

\Knowledge{
\path{durrett_ch2_joint_law_product_measure};
\path{durrett_ch2_expectation_factorization};
\path{durrett_ch2_convolution_cdf}.
}

\noindent\textbf{Solution.}
By independence, the joint law of $(X,Y)$ is $\mu\times\nu$.  The set
\[
  \{(x,y):x+y=0\}
\]
can have positive product measure only at atoms.  More explicitly, for
fixed $y$ the vertical section contains the single point $x=-y$, so
Tonelli gives
\[
  \Prob(X+Y=0)
  =\int \mu(\{-y\})\,\nu(dy).
\]
The integrand is nonzero only when $-y$ is an atom of $\mu$.  A
probability measure has at most countably many atoms, so this integral is
the countable sum
\[
  \Prob(X+Y=0)=\sum_y \mu(\{-y\})\nu(\{y\}).
\]

For (ii), apply (i) to $X$ and $-Y$.  If $X$ has continuous distribution,
then $\mu(\{y\})=0$ for every $y$, so
\[
  \Prob(X=Y)=\Prob(X-Y=0)=0.
\]

\section*{Exercise 2.1.6}

\noindent\textbf{Question.}
Prove directly from the definition that if $X$ and $Y$ are independent
and $f$ and $g$ are measurable functions, then $f(X)$ and $g(Y)$ are
independent.

\Knowledge{
\path{durrett_ch2_functions_of_independent_blocks};
\path{durrett_ch2_rv_sigma_independence}.
}

\noindent\textbf{Solution.}
Let $A,B\in\B(\R)$.  Since $f$ and $g$ are measurable,
$f^{-1}(A)$ and $g^{-1}(B)$ are Borel sets.  Therefore
\[
\begin{aligned}
  \Prob(f(X)\in A,\ g(Y)\in B)
  &=\Prob(X\in f^{-1}(A),\ Y\in g^{-1}(B))\\
  &=\Prob(X\in f^{-1}(A))\Prob(Y\in g^{-1}(B))\\
  &=\Prob(f(X)\in A)\Prob(g(Y)\in B).
\end{aligned}
\]
This is exactly the definition of independence of $f(X)$ and $g(Y)$.

\section*{Exercise 2.1.7}

\noindent\textbf{Question.}
Let $K\ge 3$ be a prime and let $X$ and $Y$ be independent random
variables uniformly distributed on $\{0,1,\ldots,K-1\}$.  For
$0\le n<K$, let
\[
  Z_n=X+nY \pmod K.
\]
Show that $Z_0,Z_1,\ldots,Z_{K-1}$ are pairwise independent.  They are
not independent because if we know the values of two of the variables
then we know the values of all the variables.

\Knowledge{
\path{durrett_ch2_event_independence};
\path{durrett_ch2_functions_of_independent_blocks};
\path{durrett_ch2_joint_law_product_measure}.
}

\noindent\textbf{Solution.}
Fix $r\ne s$ in $\{0,\ldots,K-1\}$.  For any $a,b\in\{0,\ldots,K-1\}$,
the equations
\[
  X+rY\equiv a,\qquad X+sY\equiv b \pmod K
\]
have a unique solution modulo $K$.  Indeed, subtracting gives
\[
  (s-r)Y\equiv b-a \pmod K.
\]
Since $K$ is prime and $s-r\not\equiv 0\pmod K$, $s-r$ is invertible
modulo $K$, so $Y$ is uniquely determined, and then $X$ is uniquely
determined.  Since $(X,Y)$ is uniform on $K^2$ pairs,
\[
  \Prob(Z_r=a,Z_s=b)=K^{-2}.
\]
Also each $Z_r$ is uniform on $\{0,\ldots,K-1\}$, so
\[
  \Prob(Z_r=a)\Prob(Z_s=b)=K^{-1}K^{-1}=K^{-2}.
\]
Thus $Z_r$ and $Z_s$ are independent.  Since $r\ne s$ were arbitrary,
the family is pairwise independent.

The whole family is not independent.  Any two distinct $Z_r,Z_s$
determine $X$ and $Y$, hence determine all $Z_n$.  For example,
\[
  \Prob(Z_0=0,\ldots,Z_{K-1}=0)=\Prob(X=0,Y=0)=K^{-2},
\]
whereas mutual independence would give $K^{-K}$.  Since $K\ge 3$, these
numbers are different.

\section*{Exercise 2.1.8}

\noindent\textbf{Question.}
Find four random variables taking values in $\{-1,1\}$ so that any three
are independent but all four are not.  Hint: Consider products of
independent random variables.

\Knowledge{
\path{durrett_ch2_event_independence};
\path{durrett_ch2_functions_of_independent_blocks};
\path{durrett_ch2_expectation_factorization}.
}

\noindent\textbf{Solution.}
Let $U,V,W$ be independent random variables with
\[
  \Prob(U=1)=\Prob(U=-1)=\frac12,
\]
and similarly for $V$ and $W$.  Define
\[
  X_1=U,\qquad X_2=V,\qquad X_3=W,\qquad X_4=UVW.
\]
Every $X_i$ takes values in $\{-1,1\}$.

The variables $X_1,X_2,X_3$ are independent by construction.  Consider,
for instance, $X_1,X_2,X_4$.  If $a,b,c\in\{-1,1\}$, then
\[
  \{X_1=a,X_2=b,X_4=c\}
  =\{U=a,V=b,W=abc\},
\]
which has probability $1/8$.  Since each of $X_1,X_2,X_4$ is symmetric
on $\{-1,1\}$, this equals the product of the three marginal
probabilities.  The same argument applies to every triple.

However,
\[
  X_1X_2X_3X_4=U\cdot V\cdot W\cdot UVW=1
\]
almost surely.  If all four variables were independent, then each of the
$16$ sign vectors would have probability $1/16$, including sign vectors
whose product is $-1$.  This is impossible.  Hence any three are
independent, but all four are not.

\section*{Exercise 2.1.9}

\noindent\textbf{Question.}
Let $\Omega=\{1,2,3,4\}$, $\F$ be all subsets of $\Omega$, and
$\Prob(\{i\})=1/4$.  Give an example of two collections of sets
$\mathcal A_1$ and $\mathcal A_2$ that are independent but whose
generated sigma-fields are not.

\Knowledge{
\path{durrett_ch2_pi_lambda_independence};
\path{durrett_ch2_event_independence}.
}

\noindent\textbf{Solution.}
Let
\[
  A=\{1,2\},\qquad B=\{1,3\},\qquad C=\{1,4\},
\]
and set
\[
  \mathcal A_1=\{A,B\},\qquad \mathcal A_2=\{C\}.
\]
Then
\[
  \Prob(A)=\Prob(B)=\Prob(C)=\frac12,
\]
and
\[
  \Prob(A\cap C)=\Prob(\{1\})=\frac14
  =\Prob(A)\Prob(C),
\]
while
\[
  \Prob(B\cap C)=\Prob(\{1\})=\frac14
  =\Prob(B)\Prob(C).
\]
Thus the two collections are independent.

However, $\sigma(\mathcal A_1)$ contains $\{1\}=A\cap B$, while
$C\in\sigma(\mathcal A_2)$.  These two events are not independent:
\[
  \Prob(\{1\}\cap C)=\frac14
  \ne \frac18
  =\Prob(\{1\})\Prob(C).
\]
Therefore $\sigma(\mathcal A_1)$ and $\sigma(\mathcal A_2)$ are not
independent.  This illustrates why the pi-system hypothesis in the
independence extension theorem is essential.

\section*{Exercise 2.1.10}

\noindent\textbf{Question.}
Show that if $X$ and $Y$ are independent, integer-valued random
variables, then
\[
  \Prob(X+Y=n)=\sum_m \Prob(X=m)\Prob(Y=n-m).
\]

\Knowledge{
\path{durrett_ch2_joint_law_product_measure};
\path{durrett_ch2_convolution_cdf};
\path{durrett_ch2_expectation_factorization}.
}

\noindent\textbf{Solution.}
The events $\{X=m,Y=n-m\}$, $m\in\mathbb Z$, are disjoint and their
union is $\{X+Y=n\}$.  Hence
\[
\begin{aligned}
  \Prob(X+Y=n)
  &=\sum_{m\in\mathbb Z}\Prob(X=m,Y=n-m)\\
  &=\sum_{m\in\mathbb Z}\Prob(X=m)\Prob(Y=n-m),
\end{aligned}
\]
where the last step uses independence.

\section*{Exercise 2.1.11}

\noindent\textbf{Question.}
In Example 1.6.13, the Poisson distribution with parameter $\lambda$ is
given by
\[
  \Prob(Z=k)=e^{-\lambda}\frac{\lambda^k}{k!},\qquad k=0,1,2,\ldots.
\]
Use the previous exercise to show that if $X=\operatorname{Poisson}
(\lambda)$ and $Y=\operatorname{Poisson}(\mu)$ are independent, then
$X+Y=\operatorname{Poisson}(\lambda+\mu)$.

\Knowledge{
\path{durrett_ch2_convolution_cdf};
\path{durrett_ch2_joint_law_product_measure};
\path{durrett_ch2_standard_convolution_examples}.
}

\noindent\textbf{Solution.}
For $n=0,1,2,\ldots$, Exercise 2.1.10 gives
\[
\begin{aligned}
  \Prob(X+Y=n)
  &=\sum_{m=0}^n
    e^{-\lambda}\frac{\lambda^m}{m!}
    e^{-\mu}\frac{\mu^{n-m}}{(n-m)!}\\
  &=e^{-(\lambda+\mu)}
    \frac{1}{n!}\sum_{m=0}^n \binom{n}{m}\lambda^m\mu^{n-m}\\
  &=e^{-(\lambda+\mu)}\frac{(\lambda+\mu)^n}{n!}.
\end{aligned}
\]
This is the probability mass function of a
$\operatorname{Poisson}(\lambda+\mu)$ random variable.

\section*{Exercise 2.1.12}

\noindent\textbf{Question.}
$X$ is said to have a $\operatorname{Binomial}(n,p)$ distribution if
\[
  \Prob(X=m)=\binom{n}{m}p^m(1-p)^{n-m}.
\]
\begin{enumerate}[label=(\roman*)]
\item Show that if $X=\operatorname{Binomial}(n,p)$ and
$Y=\operatorname{Binomial}(m,p)$ are independent, then
$X+Y=\operatorname{Binomial}(n+m,p)$.
\item Look at Example 1.6.12 and use induction to conclude that the sum
of $n$ independent $\operatorname{Bernoulli}(p)$ random variables is
$\operatorname{Binomial}(n,p)$.
\end{enumerate}

\Knowledge{
\path{durrett_ch2_convolution_cdf};
\path{durrett_ch2_joint_law_product_measure};
\path{durrett_ch2_standard_convolution_examples}.
}

\noindent\textbf{Solution.}
For $k=0,\ldots,n+m$, Exercise 2.1.10 gives
\[
\begin{aligned}
  \Prob(X+Y=k)
  &=\sum_j \binom{n}{j}p^j(1-p)^{n-j}
          \binom{m}{k-j}p^{k-j}(1-p)^{m-k+j}\\
  &=p^k(1-p)^{n+m-k}
    \sum_j \binom{n}{j}\binom{m}{k-j}.
\end{aligned}
\]
By Vandermonde's identity,
\[
  \sum_j \binom{n}{j}\binom{m}{k-j}=\binom{n+m}{k},
\]
so
\[
  \Prob(X+Y=k)=\binom{n+m}{k}p^k(1-p)^{n+m-k}.
\]
Therefore $X+Y$ is $\operatorname{Binomial}(n+m,p)$.

For (ii), a Bernoulli$(p)$ variable is Binomial$(1,p)$.  The case
$n=1$ is immediate.  If the sum of $n$ independent Bernoulli$(p)$
variables is Binomial$(n,p)$, then adding one more independent
Bernoulli$(p)$ variable and applying part (i) gives a
Binomial$(n+1,p)$ distribution.  The conclusion follows by induction.

\section*{Exercise 2.1.13}

\noindent\textbf{Question.}
It should not be surprising that the distribution of $X+Y$ can be
$F*G$ without the random variables being independent.  Suppose
$X,Y\in\{0,1,2\}$ and each takes each value with probability $1/3$.
\begin{enumerate}[label=(\alph*)]
\item Find the distribution of $X+Y$ assuming $X$ and $Y$ are
independent.
\item Find all the joint distributions $(X,Y)$ so that the distribution
of $X+Y$ is the same as the answer to (a).
\end{enumerate}

\Knowledge{
\path{durrett_ch2_convolution_cdf};
\path{durrett_ch2_joint_law_product_measure};
\path{durrett_ch2_event_independence}.
}

\noindent\textbf{Solution.}
If $X$ and $Y$ are independent and uniform on $\{0,1,2\}$, then the
number of pairs $(x,y)$ with $x+y=k$ is $1,2,3,2,1$ for
$k=0,1,2,3,4$.  Hence
\[
  \Prob(X+Y=k)=
  \begin{cases}
    1/9, & k=0,4,\\
    2/9, & k=1,3,\\
    3/9, & k=2.
  \end{cases}
\]

For (b), write $p_{ij}=\Prob(X=i,Y=j)$ and scale by
$q_{ij}=9p_{ij}$.  The marginal conditions say each row and each column
of $q=(q_{ij})_{0\le i,j\le 2}$ sums to $3$.  The required distribution
of $X+Y$ says the anti-diagonal sums are $1,2,3,2,1$.

Solving these linear constraints gives the one-parameter family
\[
  q=
  \begin{pmatrix}
  1 & t & 2-t\\
  2-t & 1 & t\\
  t & 2-t & 1
  \end{pmatrix},
  \qquad 0\le t\le 2.
\]
Therefore all such joint distributions are
\[
  \Prob(X=i,Y=j)=\frac{1}{9}
  \begin{pmatrix}
  1 & t & 2-t\\
  2-t & 1 & t\\
  t & 2-t & 1
  \end{pmatrix}_{ij},
  \qquad 0\le t\le 2.
\]
The independent case corresponds to $t=1$, but every $t\in[0,2]$ gives
the same distribution for $X+Y$.

\section*{Exercise 2.1.14}

\noindent\textbf{Question.}
Let $X,Y\ge 0$ be independent with distribution functions $F$ and $G$.
Find the distribution function of $XY$.

\Knowledge{
\path{durrett_ch2_joint_law_product_measure};
\path{durrett_ch2_expectation_factorization};
\path{durrett_ch2_convolution_cdf}.
}

\noindent\textbf{Solution.}
Let $H(t)=\Prob(XY\le t)$.  Since $X,Y\ge 0$, $H(t)=0$ for $t<0$.  For
$t\ge 0$, condition on the value of $Y$ and use independence:
\[
  H(t)=\int_{[0,\infty)} \Prob(Xy\le t)\,G(dy).
\]
If $y=0$, then $Xy=0\le t$, so the inner probability is $1$.  If
$y>0$, it is $\Prob(X\le t/y)=F(t/y)$.  Thus
\[
  H(t)=G(0)+\int_{(0,\infty)} F(t/y)\,G(dy),
  \qquad t\ge 0.
\]
Equivalently, with the convention $F(\infty)=1$, this can be written as
\[
  H(t)=\int_{[0,\infty)} F(t/y)\,G(dy),\qquad t\ge 0.
\]

\section*{Exercise 2.1.15}

\noindent\textbf{Question.}
If we want an infinite sequence of coin tossings, we do not have to use
Kolmogorov's theorem.  Let $\Omega$ be the unit interval $(0,1)$ equipped
with the Borel sets $\F$ and Lebesgue measure $\Prob$.  Let
\[
  Y_n(\omega)=
  \begin{cases}
  1, & \lfloor 2^n\omega\rfloor \text{ is odd},\\
  0, & \lfloor 2^n\omega\rfloor \text{ is even}.
  \end{cases}
\]
Show that $Y_1,Y_2,\ldots$ are independent with
$\Prob(Y_k=0)=\Prob(Y_k=1)=1/2$.

\Knowledge{
\path{durrett_ch2_construct_independent_variables};
\path{durrett_ch2_event_independence};
\path{durrett_ch2_pi_lambda_independence}.
}

\noindent\textbf{Solution.}
Except on the countable set of dyadic rationals, every $\omega\in(0,1)$
has a unique binary expansion
\[
  \omega=.b_1b_2b_3\cdots,\qquad b_i\in\{0,1\}.
\]
The parity of $\lfloor 2^n\omega\rfloor$ is exactly the $n$th binary
digit $b_n$.  The exceptional dyadic rationals have Lebesgue measure
zero, so they do not affect the probabilities below.

Fix $n$ and a binary vector $(\varepsilon_1,\ldots,\varepsilon_n)$.  The
set
\[
  \{\omega:Y_1=\varepsilon_1,\ldots,Y_n=\varepsilon_n\}
\]
is, up to endpoints, one dyadic interval of length $2^{-n}$: it is the
set of numbers whose first $n$ binary digits are the prescribed vector.
Therefore
\[
  \Prob(Y_1=\varepsilon_1,\ldots,Y_n=\varepsilon_n)=2^{-n}
  =\prod_{k=1}^n \frac12.
\]
Taking $n=1$ gives $\Prob(Y_k=0)=\Prob(Y_k=1)=1/2$ for each $k$, and the
finite-dimensional product formula proves that $Y_1,Y_2,\ldots$ are
independent.

\section*{Exercise 2.2.1}

\noindent\textbf{Question.}
Let $X_1,X_2,\ldots$ be uncorrelated with $\E X_i=\mu_i$ and
$\operatorname{var}(X_i)/i\to 0$ as $i\to\infty$.  Let
$S_n=X_1+\cdots+X_n$ and $\nu_n=\E S_n/n$.  Then, as $n\to\infty$,
\[
  S_n/n-\nu_n\to 0
\]
in $L^2$ and in probability.

\Knowledge{
\path{durrett_ch2_weak_law_l2};
\path{durrett_ch2_iid_finite_variance_wlln};
\path{durrett_ch2_poissonization_and_wlln_examples}.
}

\noindent\textbf{Solution.}
Since the variables are uncorrelated,
\[
  \operatorname{var}(S_n)=\sum_{i=1}^n \operatorname{var}(X_i).
\]
Thus
\[
  \E(S_n/n-\nu_n)^2
  =\operatorname{var}(S_n/n)
  =\frac{1}{n^2}\sum_{i=1}^n \operatorname{var}(X_i).
\]
Write $a_i=\operatorname{var}(X_i)/i$.  Then $a_i\to 0$ and
\[
  \frac{1}{n^2}\sum_{i=1}^n \operatorname{var}(X_i)
  =\frac{1}{n^2}\sum_{i=1}^n i a_i.
\]
Given $\varepsilon>0$, choose $N$ so that $a_i\le\varepsilon$ for
$i\ge N$.  Then
\[
  \frac{1}{n^2}\sum_{i=1}^n i a_i
  \le \frac{1}{n^2}\sum_{i=1}^{N-1} i a_i
      +\varepsilon\frac{1}{n^2}\sum_{i=N}^n i.
\]
The first term tends to $0$, and the second is at most
$\varepsilon/2+o(1)$.  Since $\varepsilon$ is arbitrary,
\[
  \E(S_n/n-\nu_n)^2\to 0.
\]
This proves convergence in $L^2$.  Convergence in probability follows
from Chebyshev's inequality.

\section*{Exercise 2.2.2}

\noindent\textbf{Question.}
The $L^2$ weak law generalizes immediately to certain dependent
sequences.  Suppose $\E X_n=0$ and
\[
  \E X_nX_m\le r(n-m),\qquad m\le n,
\]
with $r(k)\to 0$ as $k\to\infty$.  Show that
\[
  (X_1+\cdots+X_n)/n\to 0
\]
in probability.

\Knowledge{
\path{durrett_ch2_weak_law_l2};
\path{durrett_ch2_poissonization_and_wlln_examples}.
}

\noindent\textbf{Solution.}
Let $S_n=X_1+\cdots+X_n$.  Since $\E X_i=0$,
\[
  \E(S_n/n)^2=\frac{1}{n^2}\E S_n^2.
\]
Using the assumed covariance bound,
\[
\begin{aligned}
  \E S_n^2
  &=\sum_{i=1}^n \E X_i^2
    +2\sum_{1\le m<i\le n}\E X_iX_m \\
  &\le n r(0)+2\sum_{k=1}^{n-1}(n-k)r(k).
\end{aligned}
\]
Fix $\varepsilon>0$.  Choose $K$ so that $r(k)\le \varepsilon$ for
$k\ge K$.  Then
\[
  \frac{1}{n^2}\E S_n^2
  \le \frac{r(0)}{n}
      +\frac{2}{n^2}\sum_{k=1}^{K-1}(n-k)r(k)
      +\frac{2\varepsilon}{n^2}\sum_{k=K}^{n-1}(n-k).
\]
The first two terms tend to $0$, while the last term is at most
$\varepsilon+o(1)$.  Since $\E(S_n/n)^2\ge 0$ and $\varepsilon$ is
arbitrary, $\E(S_n/n)^2\to 0$.  Chebyshev's inequality gives
$S_n/n\to 0$ in probability.

\section*{Exercise 2.2.3}

\noindent\textbf{Question.}
Monte Carlo integration.
\begin{enumerate}[label=(\roman*)]
\item Let $f$ be a measurable function on $[0,1]$ with
$\int_0^1 |f(x)|\,dx<\infty$.  Let $U_1,U_2,\ldots$ be independent and
uniformly distributed on $[0,1]$, and let
\[
  I_n=n^{-1}(f(U_1)+\cdots+f(U_n)).
\]
Show that $I_n\to I\equiv\int_0^1 f\,dx$ in probability.
\item Suppose $\int_0^1 |f(x)|^2\,dx<\infty$.  Use Chebyshev's
inequality to estimate $\Prob(|I_n-I|>a/n^{1/2})$.
\end{enumerate}

\Knowledge{
\path{durrett_ch2_iid_finite_mean_wlln};
\path{durrett_ch2_iid_finite_variance_wlln};
\path{durrett_ch2_weak_law_l2}.
}

\noindent\textbf{Solution.}
Let $Y_i=f(U_i)$.  The $Y_i$ are i.i.d. and
\[
  \E |Y_1|=\int_0^1 |f(x)|\,dx<\infty,\qquad
  \E Y_1=\int_0^1 f(x)\,dx=I.
\]
By the finite-mean weak law of large numbers,
\[
  I_n=\frac{1}{n}\sum_{i=1}^n Y_i\to I
\]
in probability.

For (ii), assume $f\in L^2[0,1]$.  Then
\[
  \sigma^2=\operatorname{var}(f(U_1))
  =\int_0^1 f(x)^2\,dx-I^2<\infty.
\]
Since $\operatorname{var}(I_n)=\sigma^2/n$, Chebyshev's inequality gives
\[
  \Prob\left(|I_n-I|>\frac{a}{\sqrt n}\right)
  \le
  \frac{\sigma^2/n}{a^2/n}
  =\frac{\sigma^2}{a^2}.
\]
In particular, for a fixed error threshold $a>0$ one also has
$\Prob(|I_n-I|>a)\le \sigma^2/(na^2)$.

\section*{Exercise 2.2.4}

\noindent\textbf{Question.}
Let $X_1,X_2,\ldots$ be i.i.d. with
\[
  \Prob(X_i=(-1)^k k)=\frac{C}{k^2\log k},\qquad k\ge 2,
\]
where $C$ is chosen to make the sum of the probabilities equal to $1$.
Show that $\E|X_i|=\infty$, but there is a finite constant $\mu$ so that
$S_n/n\to\mu$ in probability.

\Knowledge{
\path{durrett_ch2_truncation_method};
\path{durrett_ch2_iid_finite_mean_wlln};
\path{durrett_ch2_infinite_mean_behavior}.
}

\noindent\textbf{Solution.}
First,
\[
  \E|X_1|
  =C\sum_{k=2}^\infty \frac{k}{k^2\log k}
  =C\sum_{k=2}^\infty \frac{1}{k\log k}
  =\infty.
\]
However,
\[
  \Prob(|X_1|>x)
  =C\sum_{k>x}\frac{1}{k^2\log k}
  \le \frac{C'}{x\log x}
\]
for large $x$, so $x\Prob(|X_1|>x)\to 0$.

By the truncation weak law, with
\[
  \mu_n=\E[X_1\mathbf 1_{\{|X_1|\le n\}}],
\]
we have
\[
  S_n/n-\mu_n\to 0
\]
in probability.  It remains only to identify the limit of $\mu_n$.
Since
\[
  \mu_n
  =C\sum_{2\le k\le n}\frac{(-1)^k}{k\log k},
\]
and the terms $1/(k\log k)$ decrease to $0$, the alternating series
test implies that $\mu_n$ converges to a finite constant
\[
  \mu=C\sum_{k=2}^\infty \frac{(-1)^k}{k\log k}.
\]
Therefore $S_n/n\to\mu$ in probability.

\section*{Exercise 2.2.5}

\noindent\textbf{Question.}
Let $X_1,X_2,\ldots$ be i.i.d. with
\[
  \Prob(X_i>x)=\frac{e}{x\log x},\qquad x\ge e.
\]
Show that $\E|X_i|=\infty$, but there is a sequence of constants
$\mu_n\to\infty$ so that
\[
  S_n/n-\mu_n\to 0
\]
in probability.

\Knowledge{
\path{durrett_ch2_truncation_method};
\path{durrett_ch2_infinite_mean_behavior};
\path{durrett_ch2_iid_finite_mean_wlln}.
}

\noindent\textbf{Solution.}
Since $X_i$ has a positive heavy tail,
\[
  \E X_1
  \ge \int_e^\infty \Prob(X_1>x)\,dx
  =\int_e^\infty \frac{e}{x\log x}\,dx
  =\infty.
\]
On the other hand,
\[
  x\Prob(|X_1|>x)=x\Prob(X_1>x)=\frac{e}{\log x}\to 0.
\]
The truncation weak law therefore applies.  Define
\[
  \mu_n=\E[X_1\mathbf 1_{\{|X_1|\le n\}}].
\]
Then
\[
  S_n/n-\mu_n\to 0
\]
in probability.  Finally, $\mu_n\to\infty$ by monotone convergence,
because $\E X_1=\infty$.

\section*{Exercise 2.2.6}

\noindent\textbf{Question.}
\begin{enumerate}[label=(\roman*)]
\item Show that if $X\ge 0$ is integer valued, then
\[
  \E X=\sum_{n\ge 1}\Prob(X\ge n).
\]
\item Find a similar expression for $\E X^2$.
\end{enumerate}

\Knowledge{
\path{durrett_ch1_tail_sum_nonnegative};
\path{durrett_ch2_truncation_method}.
}

\noindent\textbf{Solution.}
If $X=k$, then
\[
  k=\sum_{n=1}^\infty \mathbf 1_{\{n\le k\}}.
\]
Thus, for integer-valued $X\ge 0$,
\[
  X=\sum_{n=1}^\infty \mathbf 1_{\{X\ge n\}}.
\]
Taking expectations and using Tonelli's theorem gives
\[
  \E X=\sum_{n=1}^\infty \Prob(X\ge n).
\]

Similarly,
\[
  k^2=\sum_{n=1}^k (2n-1),
\]
so
\[
  X^2=\sum_{n=1}^\infty (2n-1)\mathbf 1_{\{X\ge n\}}.
\]
Therefore
\[
  \E X^2=\sum_{n=1}^\infty (2n-1)\Prob(X\ge n).
\]

\section*{Exercise 2.2.7}

\noindent\textbf{Question.}
Generalize Lemma 2.2.13 to conclude that if
\[
  H(x)=\int_{(-\infty,x]}h(y)\,dy,\qquad h(y)\ge 0,
\]
then
\[
  \E H(X)=\int_{-\infty}^\infty h(y)\Prob(X\ge y)\,dy.
\]
An important special case is $H(x)=\exp(\theta x)$ with $\theta>0$.

\Knowledge{
\path{durrett_ch1_fubini_tonelli};
\path{durrett_ch1_tail_sum_nonnegative};
\path{durrett_ch2_truncation_method}.
}

\noindent\textbf{Solution.}
Since $h\ge 0$,
\[
  H(X)=\int_{-\infty}^{\infty}h(y)\mathbf 1_{\{y\le X\}}\,dy.
\]
By Tonelli's theorem,
\[
\begin{aligned}
  \E H(X)
  &=\E\int_{-\infty}^{\infty}
    h(y)\mathbf 1_{\{y\le X\}}\,dy\\
  &=\int_{-\infty}^{\infty}
    h(y)\Prob(y\le X)\,dy\\
  &=\int_{-\infty}^{\infty}h(y)\Prob(X\ge y)\,dy.
\end{aligned}
\]
For $H(x)=e^{\theta x}$ with $\theta>0$, take
$h(y)=\theta e^{\theta y}$.  Then
\[
  \E e^{\theta X}
  =\int_{-\infty}^{\infty}\theta e^{\theta y}\Prob(X\ge y)\,dy,
\]
with both sides allowed to be $+\infty$.

\section*{Exercise 2.2.8}

\noindent\textbf{Question.}
An unfair ``fair game.''  Let
\[
  p_k=\frac{1}{2^k k(k+1)},\qquad k=1,2,\ldots,
\]
and $p_0=1-\sum_{k\ge 1}p_k$.  Since
\[
  \sum_{k=1}^\infty 2^k p_k
  =\left(1-\frac12\right)+\left(\frac12-\frac13\right)+\cdots=1,
\]
if $X_1,X_2,\ldots$ are i.i.d. with
\[
  \Prob(X_n=-1)=p_0,\qquad
  \Prob(X_n=2^k-1)=p_k,\quad k\ge 1,
\]
then $\E X_n=0$.  Let $S_n=X_1+\cdots+X_n$.  Use Theorem 2.2.11 with
$b_n=2^{m(n)}$, where
\[
  m(n)=\min\{m:2^{-m}m^{-3/2}\le n^{-1}\},
\]
to conclude that
\[
  S_n/(n/\log_2 n)\to -1
\]
in probability.

\Knowledge{
\path{durrett_ch2_triangular_array_wlln};
\path{durrett_ch2_truncation_method};
\path{durrett_ch2_infinite_mean_behavior}.
}

\noindent\textbf{Solution.}
Let $b_n=2^{m(n)}$ and truncate $X_k$ at level $b_n$:
\[
  \bar X_{n,k}=X_k\mathbf 1_{\{|X_k|\le b_n\}}.
\]
We verify the two hypotheses of Theorem 2.2.11.

First,
\[
  n\Prob(|X_1|>b_n)
  =n\sum_{j\ge m(n)+1}\frac{1}{2^j j(j+1)}
  \le C n\,2^{-m(n)}m(n)^{-2}.
\]
By the definition of $m(n)$,
$n2^{-m(n)}m(n)^{-3/2}\le 1$, so the last expression is at most
$C m(n)^{-1/2}\to 0$.

Second,
\[
\begin{aligned}
  \E \bar X_{n,1}^2
  &\le 1+\sum_{j=1}^{m(n)}
      \frac{(2^j-1)^2}{2^j j(j+1)}
   \le C\sum_{j=1}^{m(n)} \frac{2^j}{j(j+1)}
   \le C\frac{2^{m(n)}}{m(n)^2}.
\end{aligned}
\]
Hence
\[
  b_n^{-2}\sum_{k=1}^n \E \bar X_{n,k}^2
  \le C n\,2^{-m(n)}m(n)^{-2}
  \le C m(n)^{-1/2}\to 0.
\]

The centering in Theorem 2.2.11 is
\[
  a_n=n\E\bar X_{n,1}.
\]
Since $\E X_1=0$,
\[
\begin{aligned}
  \E\bar X_{n,1}
  &=-\sum_{j>m(n)}(2^j-1)p_j\\
  &=-\sum_{j>m(n)}
      \frac{1-2^{-j}}{j(j+1)}
   \sim -\frac{1}{m(n)}.
\end{aligned}
\]
Also the defining minimality of $m(n)$ implies
$m(n)/\log_2 n\to 1$ and
\[
  \frac{b_n}{n/\log_2 n}
  =\frac{2^{m(n)}\log_2 n}{n}\to 0.
\]
Therefore Theorem 2.2.11 gives
\[
  \frac{S_n-a_n}{n/\log_2 n}\to 0
\]
in probability, while
\[
  \frac{a_n}{n/\log_2 n}
  \sim -\frac{\log_2 n}{m(n)}\to -1.
\]
Combining the two limits proves
\[
  \frac{S_n}{n/\log_2 n}\to -1
\]
in probability.

\section*{Exercise 2.2.9}

\noindent\textbf{Question.}
Weak law for positive variables.  Suppose $X_1,X_2,\ldots$ are i.i.d.,
$\Prob(0\le X_i<\infty)=1$, and $\Prob(X_i>x)>0$ for all $x$.  Let
\[
  \mu(s)=\int_0^s x\,dF(x),
  \qquad
  \nu(s)=\frac{\mu(s)}{s(1-F(s))}.
\]
It is known that there exist constants $a_n$ so that
$S_n/a_n\to 1$ in probability if and only if $\nu(s)\to\infty$ as
$s\to\infty$.  Pick $b_n\ge 1$ so that $n\mu(b_n)=b_n$ for large $n$,
and use Theorem 2.2.11 to prove that the condition is sufficient.

\Knowledge{
\path{durrett_ch2_triangular_array_wlln};
\path{durrett_ch2_truncation_method};
\path{durrett_ch2_infinite_mean_behavior}.
}

\noindent\textbf{Solution.}
Assume $\nu(s)\to\infty$.  Choose $b_n$ so that
\[
  n\mu(b_n)=b_n.
\]
Since $\Prob(X_i>x)>0$ for all $x$, we have $b_n\to\infty$.

Apply Theorem 2.2.11 to the triangular array
$X_{n,k}=X_k$ with truncation level $b_n$.  The centering is
\[
  a_n=\sum_{k=1}^n \E[X_k\mathbf 1_{\{X_k\le b_n\}}]
  =n\mu(b_n)=b_n.
\]
For condition (i),
\[
  \sum_{k=1}^n \Prob(X_{n,k}>b_n)
  =n(1-F(b_n))
  =\frac{b_n(1-F(b_n))}{\mu(b_n)}
  =\frac{1}{\nu(b_n)}\to 0.
\]

For condition (ii), it is enough to show
\[
  \frac{\E[X_1^2\mathbf 1_{\{X_1\le b_n\}}]}{b_n\mu(b_n)}\to 0.
\]
Using the tail-integral identity,
\[
  \E[X_1^2\mathbf 1_{\{X_1\le b\}}]
  \le 2\int_0^b x\Prob(X_1>x)\,dx.
\]
Since $\nu(s)\to\infty$, for every $\varepsilon>0$ there is $M$ such
that
\[
  x\Prob(X_1>x)\le \varepsilon\mu(x)\le \varepsilon\mu(b),
  \qquad M\le x\le b.
\]
Thus
\[
  2\int_0^b x\Prob(X_1>x)\,dx
  \le C_M+2\varepsilon b\mu(b),
\]
where $C_M=2\int_0^M x\Prob(X_1>x)\,dx$ is fixed.  Dividing by
$b\mu(b)$ and letting $b\to\infty$ gives a limsup at most
$2\varepsilon$.  Since $\varepsilon$ is arbitrary, condition (ii) holds.

Theorem 2.2.11 now yields
\[
  \frac{S_n-b_n}{b_n}\to 0
\]
in probability.  Hence $S_n/b_n\to 1$ in probability, proving the
sufficiency of $\nu(s)\to\infty$.

\section*{Exercise 2.3.1}

\noindent\textbf{Question.}
Prove that
\[
  \Prob(\limsup A_n)\ge \limsup \Prob(A_n),
  \qquad
  \Prob(\liminf A_n)\le \liminf \Prob(A_n).
\]

\Knowledge{
\path{durrett_ch2_limsup_events};
\path{durrett_ch1_measure_continuity}.
}

\noindent\textbf{Solution.}
Let
\[
  B_m=\bigcup_{n\ge m}A_n.
\]
Then $B_m\downarrow \limsup A_n$, so continuity from above gives
\[
  \Prob(\limsup A_n)=\lim_{m\to\infty}\Prob(B_m).
\]
For each $m$, $B_m$ contains every $A_n$ with $n\ge m$, so
\[
  \Prob(B_m)\ge \sup_{n\ge m}\Prob(A_n).
\]
Letting $m\to\infty$ gives
\[
  \Prob(\limsup A_n)\ge \lim_{m\to\infty}\sup_{n\ge m}\Prob(A_n)
  =\limsup_{n\to\infty}\Prob(A_n).
\]

Similarly, let
\[
  C_m=\bigcap_{n\ge m}A_n.
\]
Then $C_m\uparrow\liminf A_n$, so
\[
  \Prob(\liminf A_n)=\lim_{m\to\infty}\Prob(C_m).
\]
Since $C_m\subseteq A_n$ for every $n\ge m$,
\[
  \Prob(C_m)\le \inf_{n\ge m}\Prob(A_n).
\]
Letting $m\to\infty$ gives
\[
  \Prob(\liminf A_n)\le \liminf_{n\to\infty}\Prob(A_n).
\]

\section*{Exercise 2.3.2}

\noindent\textbf{Question.}
Prove directly from the definition that if $f$ is continuous and
$X_n\to X$ in probability, then $f(X_n)\to f(X)$ in probability.

\Knowledge{
\path{durrett_ch2_poissonization_and_wlln_examples};
\path{durrett_ch1_composition_measurable};
\path{durrett_ch1_bounded_convergence}.
}

\noindent\textbf{Solution.}
Fix $\varepsilon>0$ and $\eta>0$.  Since $X$ is finite almost surely,
choose $M$ so large that
\[
  \Prob(|X|>M)<\eta.
\]
The function $f$ is uniformly continuous on the compact interval
$[-M-1,M+1]$.  Hence there is $0<\delta<1$ such that if
$x,y\in[-M-1,M+1]$ and $|x-y|<\delta$, then
$|f(x)-f(y)|<\varepsilon$.

On the event $\{|X|\le M,\ |X_n-X|<\delta\}$, we have
$X_n\in[-M-1,M+1]$, so $|f(X_n)-f(X)|<\varepsilon$.  Therefore
\[
  \Prob(|f(X_n)-f(X)|\ge\varepsilon)
  \le \Prob(|X|>M)+\Prob(|X_n-X|\ge\delta).
\]
Taking $\limsup$ and using $X_n\to X$ in probability,
\[
  \limsup_{n\to\infty}\Prob(|f(X_n)-f(X)|\ge\varepsilon)\le \eta.
\]
Since $\eta>0$ is arbitrary, the left-hand side is $0$.  Thus
$f(X_n)\to f(X)$ in probability.

\section*{Exercise 2.3.3}

\noindent\textbf{Question.}
Let $\ell_n$ be the length of the head run at time $n$.  Show that
\[
  \limsup_{n\to\infty}\frac{\ell_n}{\log_2 n}=1,
  \qquad
  \liminf_{n\to\infty}\ell_n=0
\]
almost surely.

\Knowledge{
\path{durrett_ch2_borel_cantelli_first};
\path{durrett_ch2_borel_cantelli_second};
\path{durrett_ch2_bc_independent_trials}.
}

\noindent\textbf{Solution.}
Assume the tosses are i.i.d. with
$\Prob(\text{head})=\Prob(\text{tail})=1/2$.  For the upper bound, fix
$\varepsilon>0$.  Since
\[
  \Prob(\ell_n\ge (1+\varepsilon)\log_2 n)
  \le 2^{-(1+\varepsilon)\log_2 n}=n^{-(1+\varepsilon)},
\]
the series of these probabilities is finite.  By the first
Borel-Cantelli lemma,
\[
  \ell_n<(1+\varepsilon)\log_2 n
\]
eventually almost surely.  Hence
\[
  \limsup_{n\to\infty}\frac{\ell_n}{\log_2 n}\le 1
  \qquad\text{a.s.}
\]

For the lower bound, fix $\varepsilon\in(0,1)$ and look at the dyadic
block of times $\{2^k+1,\ldots,2^{k+1}\}$.  Split this block into
disjoint subblocks of length
\[
  r_k=\lfloor(1-\varepsilon)k\rfloor.
\]
The probability that a given subblock consists entirely of heads is
$2^{-r_k}\ge 2^{-(1-\varepsilon)k-1}$.  There are at least
$c2^k/k$ disjoint subblocks, so the probability that none of them is all
heads is bounded by
\[
  \left(1-2^{-r_k}\right)^{c2^k/k}
  \le \exp\{-c'2^{\varepsilon k}/k\}.
\]
This is summable in $k$.  By Borel-Cantelli, eventually every dyadic
block contains a run of heads of length at least $r_k$.  Therefore
\[
  \limsup_{n\to\infty}\frac{\ell_n}{\log_2 n}\ge 1-\varepsilon
  \qquad\text{a.s.}
\]
Letting $\varepsilon\downarrow0$ gives the lower bound $1$.

Finally, $\ell_n=0$ whenever the $n$th toss is a tail.  The tail events
are independent and have probabilities summing to infinity, so by the
second Borel-Cantelli lemma tails occur infinitely often.  Hence
\[
  \liminf_{n\to\infty}\ell_n=0
  \qquad\text{a.s.}
\]

\section*{Exercise 2.3.4}

\noindent\textbf{Question.}
Fatou's lemma.  Suppose $X_n\ge 0$ and $X_n\to X$ in probability.  Show
that
\[
  \liminf_{n\to\infty}\E X_n\ge \E X.
\]

\Knowledge{
\path{durrett_ch1_fatou_integral};
\path{durrett_ch2_subsequence_bc_upgrade};
\path{durrett_ch2_limsup_events}.
}

\noindent\textbf{Solution.}
Let $a=\liminf_n \E X_n$.  If $a=\infty$, there is nothing to prove.
Otherwise choose a subsequence $n_k$ such that
\[
  \E X_{n_k}\to a.
\]
Since $X_n\to X$ in probability, the subsequence principle gives a
further subsequence $X_{n_{k_j}}$ such that
\[
  X_{n_{k_j}}\to X\qquad\text{a.s.}
\]
By the ordinary Fatou lemma for almost sure convergence and nonnegative
random variables,
\[
  \E X
  \le \liminf_{j\to\infty}\E X_{n_{k_j}}
  =a.
\]
Thus $\E X\le\liminf_n\E X_n$.

\section*{Exercise 2.3.5}

\noindent\textbf{Question.}
Dominated convergence.  Suppose $X_n\to X$ in probability and either
\begin{enumerate}[label=(\alph*)]
\item $|X_n|\le Y$ with $\E Y<\infty$, or
\item there is a continuous function $g$ with $g(x)>0$ for large $x$ and
$|x|/g(x)\to0$ as $|x|\to\infty$ so that $\E g(X_n)\le C<\infty$ for
all $n$.
\end{enumerate}
Show that $\E X_n\to \E X$.

\Knowledge{
\path{durrett_ch1_dominated_convergence};
\path{durrett_ch1_bounded_convergence};
\path{durrett_ch2_subsequence_bc_upgrade}.
}

\noindent\textbf{Solution.}
For (a), every subsequence of $X_n$ has a further subsequence
$X_{n_{k_j}}$ that converges to $X$ almost surely.  Since
$|X_{n_{k_j}}|\le Y$ and $Y$ is integrable, dominated convergence gives
\[
  \E X_{n_{k_j}}\to \E X.
\]
Thus every subsequence of the numerical sequence $\E X_n$ has a further
subsequence converging to $\E X$, so $\E X_n\to\E X$.

For (b), add a constant to $g$ if necessary so that $g\ge 1$ everywhere;
this only changes the constant $C$.  Since $g$ is continuous and
$X_n\to X$ in probability, Exercise 2.3.2 gives
$g(X_n)\to g(X)$ in probability.  By Exercise 2.3.4,
\[
  \E g(X)\le \liminf_n \E g(X_n)\le C.
\]
We now prove uniform integrability of $X_n$.  Given $\varepsilon>0$,
choose $K$ so large that
\[
  |x|\le \varepsilon g(x),\qquad |x|>K.
\]
Then, for all $n$,
\[
  \E\bigl(|X_n|\mathbf 1_{\{|X_n|>K\}}\bigr)
  \le \varepsilon \E g(X_n)
  \le \varepsilon C,
\]
and the same argument gives
\[
  \E\bigl(|X|\mathbf 1_{\{|X|>K\}}\bigr)
  \le \varepsilon C.
\]
Let $\phi_K(x)=(-K)\vee(x\wedge K)$.  This function is bounded and
continuous, so $\E\phi_K(X_n)\to\E\phi_K(X)$.  Therefore
\[
  \limsup_{n\to\infty}|\E X_n-\E X|
  \le 2\varepsilon C.
\]
Since $\varepsilon$ is arbitrary, $\E X_n\to\E X$.

\section*{Exercise 2.3.6}

\noindent\textbf{Question.}
Metric for convergence in probability.  Show
\begin{enumerate}[label=(\alph*)]
\item that
\[
  d(X,Y)=\E\left(\frac{|X-Y|}{1+|X-Y|}\right)
\]
defines a metric on the set of random variables, i.e.
$d(X,Y)=0$ if and only if $X=Y$ a.s., $d(X,Y)=d(Y,X)$, and
$d(X,Z)\le d(X,Y)+d(Y,Z)$;
\item that $d(X_n,X)\to0$ if and only if $X_n\to X$ in probability.
\end{enumerate}

\Knowledge{
\path{durrett_ch2_poissonization_and_wlln_examples};
\path{durrett_ch1_markov_chebyshev};
\path{durrett_ch1_bounded_convergence}.
}

\noindent\textbf{Solution.}
Let $\varphi(t)=t/(1+t)$ for $t\ge0$.  Then $\varphi(t)\ge0$ and
$\varphi(t)=0$ if and only if $t=0$.  Hence $d(X,Y)=0$ if and only if
$|X-Y|=0$ a.s., i.e. $X=Y$ a.s.  Symmetry is immediate.

For the triangle inequality, note that $\varphi$ is increasing and
subadditive:
\[
  \varphi(a+b)\le \varphi(a)+\varphi(b),\qquad a,b\ge0.
\]
This follows, for example, from the metric transform proved in Exercise
2.1.3.  Since
$|X-Z|\le |X-Y|+|Y-Z|$,
\[
  \varphi(|X-Z|)
  \le \varphi(|X-Y|)+\varphi(|Y-Z|).
\]
Taking expectations gives $d(X,Z)\le d(X,Y)+d(Y,Z)$.

Now suppose $X_n\to X$ in probability.  Then
$\varphi(|X_n-X|)\to0$ in probability and
$0\le \varphi(|X_n-X|)\le1$.  For any $\eta>0$,
\[
  \E\varphi(|X_n-X|)
  \le \eta+\Prob(\varphi(|X_n-X|)>\eta),
\]
and the right side has limsup at most $\eta$.  Letting $\eta\downarrow0$
shows $d(X_n,X)\to0$.

Conversely, if $d(X_n,X)\to0$, then for every $\varepsilon>0$,
\[
  \Prob(|X_n-X|>\varepsilon)
  =\Prob\left(\varphi(|X_n-X|)
       >\frac{\varepsilon}{1+\varepsilon}\right)
  \le \frac{1+\varepsilon}{\varepsilon}d(X_n,X),
\]
which tends to $0$.  Thus $X_n\to X$ in probability.

\section*{Exercise 2.3.7}

\noindent\textbf{Question.}
Show that random variables are a complete space under the metric defined
in the previous exercise, i.e. if $d(X_m,X_n)\to0$ whenever
$m,n\to\infty$, then there is a random variable $X_\infty$ so that
$X_n\to X_\infty$ in probability.

\Knowledge{
\path{durrett_ch2_borel_cantelli_first};
\path{durrett_ch2_limsup_events};
\path{durrett_ch2_subsequence_bc_upgrade}.
}

\noindent\textbf{Solution.}
By Exercise 2.3.6, $d$-Cauchy is the same as Cauchy in probability.
Choose a subsequence $n_k$ so that
\[
  \Prob(|X_{n_{k+1}}-X_{n_k}|>2^{-k})\le 2^{-k}.
\]
By Borel-Cantelli, with probability one,
\[
  |X_{n_{k+1}}-X_{n_k}|\le 2^{-k}
\]
for all sufficiently large $k$.  Hence $X_{n_k}$ is almost surely a
Cauchy sequence in $\R$ and therefore converges almost surely to some
finite random variable $X_\infty$.

It remains to show that the whole sequence converges to $X_\infty$ in
probability.  Fix $\varepsilon>0$ and $\eta>0$.  Since $X_n$ is Cauchy
in probability, choose $K$ so large that for all $m,n\ge n_K$,
\[
  \Prob(|X_m-X_n|>\varepsilon/2)<\eta.
\]
Also choose $K$ larger if necessary so that
\[
  \Prob(|X_{n_K}-X_\infty|>\varepsilon/2)<\eta.
\]
Then for all $n\ge n_K$,
\[
  \Prob(|X_n-X_\infty|>\varepsilon)
  \le \Prob(|X_n-X_{n_K}|>\varepsilon/2)
      +\Prob(|X_{n_K}-X_\infty|>\varepsilon/2)
  <2\eta.
\]
Since $\eta$ is arbitrary, $X_n\to X_\infty$ in probability.

\section*{Exercise 2.3.8}

\noindent\textbf{Question.}
Let $A_n$ be a sequence of independent events with $\Prob(A_n)<1$ for
all $n$.  Show that
\[
  \Prob\left(\bigcup_n A_n\right)=1
\]
implies $\sum_n\Prob(A_n)=\infty$ and hence
\[
  \Prob(A_n\ \text{i.o.})=1.
\]

\Knowledge{
\path{durrett_ch2_borel_cantelli_second};
\path{durrett_ch2_bc_independent_trials};
\path{durrett_ch2_event_independence}.
}

\noindent\textbf{Solution.}
Suppose, toward a contradiction, that
\[
  \sum_{n=1}^\infty \Prob(A_n)<\infty.
\]
Since $\Prob(A_n)<1$ for every $n$ and $\Prob(A_n)\to0$, the infinite
product
\[
  \prod_{n=1}^\infty (1-\Prob(A_n))
\]
is strictly positive.  Indeed, after finitely many terms we have
$\Prob(A_n)\le1/2$, and then
\[
  \log(1-\Prob(A_n))\ge -2\Prob(A_n),
\]
so the tail product is positive.

By independence,
\[
  \Prob\left(\bigcap_{n=1}^\infty A_n^c\right)
  =\prod_{n=1}^\infty (1-\Prob(A_n))>0.
\]
This contradicts $\Prob(\cup_n A_n)=1$.  Hence
$\sum_n\Prob(A_n)=\infty$.  The second Borel-Cantelli lemma then gives
\[
  \Prob(A_n\ \text{i.o.})=1.
\]

\section*{Exercise 2.3.9}

\noindent\textbf{Question.}
\begin{enumerate}[label=(\roman*)]
\item If $\Prob(A_n)\to0$ and
\[
  \sum_{n=1}^\infty \Prob(A_n^c\cap A_{n+1})<\infty,
\]
then $\Prob(A_n\ \text{i.o.})=0$.
\item Find an example of a sequence $A_n$ to which the result in (i) can
be applied but the Borel-Cantelli lemma cannot.
\end{enumerate}

\Knowledge{
\path{durrett_ch2_borel_cantelli_first};
\path{durrett_ch2_limsup_events};
\path{durrett_ch2_bc_tail_bounds}.
}

\noindent\textbf{Solution.}
Let $B_n=A_n^c\cap A_{n+1}$.  If $A_n$ occurs infinitely often, then
either $A_n$ occurs eventually always, or there are infinitely many
transitions from $A_n^c$ to $A_{n+1}$.  Thus
\[
  \limsup A_n
  \subseteq \liminf A_n\ \cup\ \limsup B_n.
\]
By Exercise 2.3.1,
\[
  \Prob(\liminf A_n)\le \liminf_n \Prob(A_n)=0.
\]
Also, since $\sum_n\Prob(B_n)<\infty$, the first Borel-Cantelli lemma
gives
\[
  \Prob(\limsup B_n)=0.
\]
Therefore $\Prob(A_n\ \text{i.o.})=0$.

For (ii), take $\Omega=(0,1)$ with Lebesgue measure and
\[
  A_n=(0,1/n).
\]
Then $\Prob(A_n)=1/n$, so $\sum_n\Prob(A_n)=\infty$ and the ordinary
Borel-Cantelli lemma cannot be applied.  However $A_{n+1}\subseteq A_n$,
so
\[
  A_n^c\cap A_{n+1}=\varnothing
\]
for every $n$, and $\Prob(A_n)\to0$.  Part (i) applies and gives
$\Prob(A_n\ \text{i.o.})=0$.

\section*{Exercise 2.3.10}

\noindent\textbf{Question.}
Kochen-Stone lemma.  Suppose $\sum \Prob(A_k)=\infty$.  Use Exercises
1.6.6 and 2.3.1 to show that if
\[
  \limsup_{n\to\infty}
  \frac{\left(\sum_{k=1}^n \Prob(A_k)\right)^2}
       {\sum_{1\le j,k\le n}\Prob(A_j\cap A_k)}
  =\alpha>0,
\]
then
\[
  \Prob(A_n\ \text{i.o.})\ge \alpha.
\]
The case $\alpha=1$ contains Theorem 2.3.7.

\Knowledge{
\path{durrett_ch2_limsup_events};
\path{durrett_ch2_borel_cantelli_second};
\path{durrett_ch2_indicator_independence};
\path{durrett_ch1_markov_chebyshev}.
}

\noindent\textbf{Solution.}
Let
\[
  S_n=\sum_{k=1}^n \mathbf 1_{A_k},
  \qquad
  m_n=\E S_n=\sum_{k=1}^n\Prob(A_k).
\]
Then $m_n\to\infty$ and
\[
  \E S_n^2=\sum_{1\le j,k\le n}\Prob(A_j\cap A_k).
\]
Exercise 1.6.6, applied to the nonnegative random variable $S_n$, gives
the second-moment lower bound
\[
  \Prob(S_n>0)\ge \frac{(\E S_n)^2}{\E S_n^2}
  =\frac{m_n^2}{\sum_{1\le j,k\le n}\Prob(A_j\cap A_k)}.
\]
Thus
\[
  \Prob\left(\bigcup_{k=1}^n A_k\right)
  \ge
  \frac{m_n^2}{\sum_{1\le j,k\le n}\Prob(A_j\cap A_k)}.
\]

We need the same estimate for tails.  Fix $r\ge1$ and define
\[
  S_{r,n}=\sum_{k=r}^n\mathbf 1_{A_k},\qquad
  m_{r,n}=\sum_{k=r}^n\Prob(A_k).
\]
The same second-moment argument gives
\[
  \Prob\left(\bigcup_{k=r}^n A_k\right)
  \ge
  \frac{m_{r,n}^2}
       {\sum_{r\le j,k\le n}\Prob(A_j\cap A_k)}.
\]
Along a subsequence for which the original ratio tends to $\alpha$, the
finite number of omitted initial terms and cross terms are lower order:
the numerator changes by $o(m_n^2)$, and the denominator changes by at
most $O(m_n)$ plus a constant.  Hence
\[
  \limsup_{n\to\infty}
  \frac{m_{r,n}^2}
       {\sum_{r\le j,k\le n}\Prob(A_j\cap A_k)}
  \ge \alpha.
\]
Therefore
\[
  \Prob\left(\bigcup_{k=r}^\infty A_k\right)\ge \alpha
  \qquad\text{for every }r.
\]
Finally,
\[
  \limsup A_n=\bigcap_{r=1}^\infty \bigcup_{k=r}^\infty A_k,
\]
and the sets $\cup_{k=r}^\infty A_k$ decrease as $r$ increases.  By
continuity from above,
\[
  \Prob(A_n\ \text{i.o.})
  =\lim_{r\to\infty}\Prob\left(\bigcup_{k=r}^\infty A_k\right)
  \ge \alpha.
\]

\section*{Exercise 2.3.11}

\noindent\textbf{Question.}
Let $X_1,X_2,\ldots$ be independent with
$\Prob(X_n=1)=p_n$ and $\Prob(X_n=0)=1-p_n$.  Show that
\begin{enumerate}[label=(\roman*)]
\item $X_n\to0$ in probability if and only if $p_n\to0$;
\item $X_n\to0$ almost surely if and only if $\sum_n p_n<\infty$.
\end{enumerate}

\Knowledge{
\path{durrett_ch2_borel_cantelli_first};
\path{durrett_ch2_borel_cantelli_second};
\path{durrett_ch2_limsup_events}.
}

\noindent\textbf{Solution.}
For $0<\varepsilon<1$,
\[
  \Prob(|X_n|>\varepsilon)=\Prob(X_n=1)=p_n.
\]
Thus $X_n\to0$ in probability if and only if $p_n\to0$.

For almost sure convergence, observe that $X_n\to0$ a.s. if and only if
the events $\{X_n=1\}$ occur only finitely often.  If
$\sum_n p_n<\infty$, the first Borel-Cantelli lemma gives
\[
  \Prob(X_n=1\ \text{i.o.})=0.
\]
Conversely, if $\sum_n p_n=\infty$, then the events $\{X_n=1\}$ are
independent, so the second Borel-Cantelli lemma gives
\[
  \Prob(X_n=1\ \text{i.o.})=1.
\]
Hence $X_n\not\to0$ a.s. in that case.

\section*{Exercise 2.3.12}

\noindent\textbf{Question.}
Let $X_1,X_2,\ldots$ be a sequence of random variables on
$(\Omega,\F,\Prob)$, where $\Omega$ is a countable set and $\F$ consists
of all subsets of $\Omega$.  Show that $X_n\to X$ in probability implies
$X_n\to X$ almost surely.

\Knowledge{
\path{durrett_ch2_limsup_events};
\path{durrett_ch2_bc_tail_bounds}.
}

\noindent\textbf{Solution.}
Let
\[
  \Omega_+=\{\omega\in\Omega:\Prob(\{\omega\})>0\}.
\]
Since $\Omega$ is countable, $\Prob(\Omega_+)=1$.  Fix
$\omega\in\Omega_+$.  We claim that $X_n(\omega)\to X(\omega)$.

If not, then for some $\varepsilon>0$ there is a subsequence $n_k$ such
that
\[
  |X_{n_k}(\omega)-X(\omega)|>\varepsilon
  \qquad\text{for all }k.
\]
Hence
\[
  \Prob(|X_{n_k}-X|>\varepsilon)\ge \Prob(\{\omega\})>0
\]
for all $k$, contradicting convergence in probability.  Therefore
$X_n(\omega)\to X(\omega)$ for every $\omega\in\Omega_+$, and so
$X_n\to X$ almost surely.

\section*{Exercise 2.3.13}

\noindent\textbf{Question.}
If $X_n$ is any sequence of random variables, there are constants
$c_n\to\infty$ so that
\[
  X_n/c_n\to0
\]
almost surely.

\Knowledge{
\path{durrett_ch2_borel_cantelli_first};
\path{durrett_ch2_bc_tail_bounds}.
}

\noindent\textbf{Solution.}
For each $n$, choose $c_n\ge n$ so large that
\[
  \Prob(|X_n|>c_n/j)\le 2^{-n},
  \qquad j=1,\ldots,n.
\]
This is possible because $\Prob(|X_n|>t)\to0$ as $t\to\infty$ for each
fixed $n$.  Clearly $c_n\to\infty$.

Fix $\varepsilon>0$ and choose an integer $j$ with $1/j<\varepsilon$.
For all $n\ge j$,
\[
  \Prob(|X_n|>\varepsilon c_n)
  \le \Prob(|X_n|>c_n/j)\le 2^{-n}.
\]
Thus
\[
  \sum_{n=1}^\infty \Prob(|X_n|>\varepsilon c_n)<\infty.
\]
By Borel-Cantelli, $|X_n|>\varepsilon c_n$ only finitely often almost
surely.  Applying this to $\varepsilon=1/m$, $m=1,2,\ldots$, gives
$X_n/c_n\to0$ almost surely.

\section*{Exercise 2.3.14}

\noindent\textbf{Question.}
Let $X_1,X_2,\ldots$ be independent.  Show that
\[
  \sup_n X_n<\infty\quad\text{a.s.}
\]
if and only if
\[
  \sum_n \Prob(X_n>A)<\infty
\]
for some $A$.

\Knowledge{
\path{durrett_ch2_borel_cantelli_first};
\path{durrett_ch2_borel_cantelli_second};
\path{durrett_ch2_bc_tail_bounds}.
}

\noindent\textbf{Solution.}
If $\sum_n\Prob(X_n>A)<\infty$, then by the first Borel-Cantelli lemma,
$X_n>A$ only finitely often a.s.  Thus all sufficiently large $n$ have
$X_n\le A$.  The maximum of the remaining finitely many finite random
variables is finite a.s., so $\sup_n X_n<\infty$ a.s.

Conversely, suppose $\sup_n X_n<\infty$ a.s.  If
\[
  \sum_n\Prob(X_n>A)=\infty
\]
for every integer $A$, then, by independence and the second
Borel-Cantelli lemma,
\[
  \Prob(X_n>A\ \text{i.o.})=1
  \qquad\text{for every integer }A.
\]
Intersecting over $A=1,2,\ldots$ gives $\sup_n X_n=\infty$ a.s., a
contradiction.  Hence for some integer $A$ the series is finite.

\section*{Exercise 2.3.15}

\noindent\textbf{Question.}
Let $X_1,X_2,\ldots$ be i.i.d. with
\[
  \Prob(X_i>x)=e^{-x},
\]
and let $M_n=\max_{1\le m\le n}X_m$.  Show that
\begin{enumerate}[label=(\roman*)]
\item $\limsup_{n\to\infty}X_n/\log n=1$ a.s.;
\item $M_n/\log n\to1$ a.s.
\end{enumerate}

\Knowledge{
\path{durrett_ch2_borel_cantelli_first};
\path{durrett_ch2_borel_cantelli_second};
\path{durrett_ch2_bc_independent_trials}.
}

\noindent\textbf{Solution.}
Fix $\varepsilon>0$.  Since
\[
  \Prob(X_n>(1+\varepsilon)\log n)=n^{-(1+\varepsilon)},
\]
the first Borel-Cantelli lemma implies
\[
  X_n\le (1+\varepsilon)\log n
\]
eventually a.s.  Hence
\[
  \limsup_{n\to\infty}\frac{X_n}{\log n}\le 1+\varepsilon
  \qquad\text{a.s.}
\]
Letting $\varepsilon\downarrow0$ gives the upper bound.

For the lower bound,
\[
  \Prob(X_n>(1-\varepsilon)\log n)=n^{-(1-\varepsilon)},
\]
and the series over $n$ diverges.  The events are independent, so by the
second Borel-Cantelli lemma they occur infinitely often.  Hence
\[
  \limsup_{n\to\infty}\frac{X_n}{\log n}\ge 1-\varepsilon
  \qquad\text{a.s.}
\]
Letting $\varepsilon\downarrow0$ proves (i).

For (ii), the upper bound in (i) implies that eventually
$X_n\le(1+\varepsilon)\log n$, and hence
\[
  M_n\le C(\omega)\vee(1+\varepsilon)\log n,
\]
so $\limsup M_n/\log n\le1+\varepsilon$ a.s.
For the lower bound,
\[
  \Prob(M_n\le(1-\varepsilon)\log n)
  =\left(1-n^{-(1-\varepsilon)}\right)^n
  \le \exp(-n^\varepsilon).
\]
The last bound is summable, so by Borel-Cantelli
$M_n>(1-\varepsilon)\log n$ eventually a.s.  Letting
$\varepsilon\downarrow0$ gives $M_n/\log n\to1$ a.s.

\section*{Exercise 2.3.16}

\noindent\textbf{Question.}
Let $X_1,X_2,\ldots$ be i.i.d. with distribution $F$, let
$\lambda_n\uparrow\infty$, and let
\[
  A_n=\left\{\max_{1\le m\le n}X_m>\lambda_n\right\}.
\]
Show that $\Prob(A_n\ \text{i.o.})=0$ or $1$ according as
\[
  \sum_{n\ge1}(1-F(\lambda_n))<\infty
  \quad\text{or}\quad
  =\infty.
\]

\Knowledge{
\path{durrett_ch2_borel_cantelli_first};
\path{durrett_ch2_borel_cantelli_second};
\path{durrett_ch2_bc_tail_bounds}.
}

\noindent\textbf{Solution.}
Let $B_n=\{X_n>\lambda_n\}$.  Then
\[
  \Prob(B_n)=1-F(\lambda_n).
\]
If $\sum_n\Prob(B_n)<\infty$, then $B_n$ occurs only finitely often
a.s.  If $A_n$ occurs infinitely often, then infinitely often some
$m\le n$ has $X_m>\lambda_n$.  Since $\lambda_n\to\infty$, a fixed
$X_m$ can only be responsible for finitely many such $n$.  Therefore
$A_n$ i.o. would force $B_m$ i.o., which has probability $0$.  Hence
$\Prob(A_n\ \text{i.o.})=0$.

If $\sum_n\Prob(B_n)=\infty$, then the independent events $B_n$ occur
infinitely often by the second Borel-Cantelli lemma.  Since
$B_n\subseteq A_n$, it follows that $A_n$ occurs infinitely often with
probability $1$.

\section*{Exercise 2.3.17}

\noindent\textbf{Question.}
Let $Y_1,Y_2,\ldots$ be i.i.d.  Find necessary and sufficient conditions
for
\begin{enumerate}[label=(\roman*)]
\item $Y_n/n\to0$ almost surely;
\item $(\max_{m\le n}Y_m)/n\to0$ almost surely;
\item $(\max_{m\le n}Y_m)/n\to0$ in probability;
\item $Y_n/n\to0$ in probability.
\end{enumerate}

\Knowledge{
\path{durrett_ch2_borel_cantelli_first};
\path{durrett_ch2_borel_cantelli_second};
\path{durrett_ch2_max_negligible_condition}.
}

\noindent\textbf{Solution.}
For (i), $Y_n/n\to0$ a.s. if and only if for every $\varepsilon>0$,
\[
  \Prob(|Y_n|>\varepsilon n\ \text{i.o.})=0.
\]
By Borel-Cantelli and independence, this is equivalent to
\[
  \sum_{n=1}^\infty \Prob(|Y_1|>\varepsilon n)<\infty
  \qquad\text{for every }\varepsilon>0.
\]
By the tail-sum criterion, this is equivalent to
\[
  \E |Y_1|<\infty.
\]

For (ii), since $\max_{m\le n}Y_m\ge Y_1$ and $Y_1/n\to0$, only the
upper tail matters.  The condition is
\[
  \sum_{n=1}^\infty \Prob(Y_1>\varepsilon n)<\infty
  \qquad\text{for every }\varepsilon>0,
\]
which is equivalent to
\[
  \E(Y_1^+)<\infty.
\]

For (iii),
\[
  \Prob\left(\max_{m\le n}Y_m>\varepsilon n\right)
  =1-\left(1-\Prob(Y_1>\varepsilon n)\right)^n.
\]
This tends to $0$ if and only if
\[
  n\Prob(Y_1>\varepsilon n)\to0
  \qquad\text{for every }\varepsilon>0,
\]
equivalently
\[
  x\Prob(Y_1>x)\to0\qquad (x\to\infty).
\]

For (iv), for every $\varepsilon>0$,
\[
  \Prob(|Y_n|>\varepsilon n)=\Prob(|Y_1|>\varepsilon n)\to0,
\]
provided $Y_1$ is finite a.s.  Thus no moment condition is needed beyond
finiteness of the random variable.

\section*{Exercise 2.3.18}

\noindent\textbf{Question.}
Let $0\le X_1\le X_2\le\cdots$ be random variables with
\[
  \E X_n\sim a n^\alpha,\qquad a,\alpha>0,
\]
and
\[
  \operatorname{var}(X_n)\le Bn^\beta,\qquad \beta<2\alpha.
\]
Show that
\[
  X_n/n^\alpha\to a
\]
almost surely.

\Knowledge{
\path{durrett_ch2_subsequence_bc_upgrade};
\path{durrett_ch1_markov_chebyshev};
\path{durrett_ch2_slln_normalized_sums}.
}

\noindent\textbf{Solution.}
Choose $r>0$ so large that
\[
  r(2\alpha-\beta)>1,
\]
and let $n_k=\lfloor k^r\rfloor$.  By Chebyshev's inequality,
\[
\begin{aligned}
  \Prob\left(\left|\frac{X_{n_k}-\E X_{n_k}}{n_k^\alpha}\right|
  >\varepsilon\right)
  &\le \frac{\operatorname{var}(X_{n_k})}
           {\varepsilon^2 n_k^{2\alpha}}\\
  &\le C k^{-r(2\alpha-\beta)}.
\end{aligned}
\]
The last series is summable.  By Borel-Cantelli,
\[
  \frac{X_{n_k}-\E X_{n_k}}{n_k^\alpha}\to0
  \qquad\text{a.s.}
\]
Since $\E X_n/n^\alpha\to a$,
\[
  X_{n_k}/n_k^\alpha\to a
  \qquad\text{a.s.}
\]

Now let $n_k\le n<n_{k+1}$.  By monotonicity,
\[
  \frac{X_{n_k}}{n_{k+1}^\alpha}
  \le \frac{X_n}{n^\alpha}
  \le \frac{X_{n_{k+1}}}{n_k^\alpha}.
\]
Because $n_{k+1}/n_k\to1$, both outside terms converge to $a$ a.s.
Therefore $X_n/n^\alpha\to a$ a.s.

\section*{Exercise 2.3.19}

\noindent\textbf{Question.}
Let $X_n$ be independent Poisson random variables with
\[
  \E X_n=\lambda_n,
\]
and let $S_n=X_1+\cdots+X_n$.  Show that if
$\sum_n\lambda_n=\infty$, then
\[
  S_n/\E S_n\to1
\]
almost surely.

\Knowledge{
\path{durrett_ch2_kolmogorov_convergence_criterion};
\path{durrett_ch2_slln_normalized_sums};
\path{durrett_ch2_weak_law_l2}.
}

\noindent\textbf{Solution.}
Let
\[
  m_n=\E S_n=\sum_{k=1}^n\lambda_k.
\]
Then $m_n\to\infty$.  Since the $X_k$ are independent Poisson random
variables,
\[
  \operatorname{var}(S_n)=\sum_{k=1}^n\operatorname{var}(X_k)
  =\sum_{k=1}^n\lambda_k=m_n.
\]
It is enough to prove
\[
  \frac{S_n-m_n}{m_n}\to0\qquad\text{a.s.}
\]
For all sufficiently large $n$, $m_n>0$.  Put
\[
  Z_n=\frac{X_n-\lambda_n}{m_n}.
\]
Then the $Z_n$ are independent, mean zero, and
\[
  \sum_n \operatorname{var}(Z_n)
  =\sum_n\frac{\lambda_n}{m_n^2}<\infty.
\]
The last series is finite because $\lambda_n=m_n-m_{n-1}$ and, once
$m_{n-1}>0$,
\[
  \frac{m_n-m_{n-1}}{m_n^2}
  \le \int_{m_{n-1}}^{m_n}\frac{dx}{x^2}.
\]
Thus $\sum_n Z_n$ converges almost surely by the Kolmogorov convergence
criterion.  Kronecker's lemma with $a_n=m_n$ then gives
\[
  \frac{1}{m_n}\sum_{k=1}^n (X_k-\lambda_k)\to0
  \qquad\text{a.s.}
\]
This is exactly $S_n/m_n\to1$ a.s.

\section*{Exercise 2.3.20}

\noindent\textbf{Question.}
Show that if $X_n$ is the outcome of the $n$th play of the St.
Petersburg game, then
\[
  \limsup_{n\to\infty}\frac{X_n}{n\log_2 n}=\infty
  \qquad\text{a.s.}
\]
and hence the same result holds for $S_n=X_1+\cdots+X_n$.  This shows
that the convergence $S_n/(n\log_2 n)\to1$ in probability proved in
Section 2.2 does not occur a.s.

\Knowledge{
\path{durrett_ch2_borel_cantelli_second};
\path{durrett_ch2_infinite_mean_behavior};
\path{durrett_ch2_unfair_fair_game_truncation_scale}.
}

\noindent\textbf{Solution.}
In the St. Petersburg game,
\[
  \Prob(X_n=2^j)=2^{-j},\qquad j\ge1.
\]
There is a constant $c>0$ such that for all large $x$,
\[
  \Prob(X_n\ge x)\ge \frac{c}{x}.
\]
Fix $M>0$ and define
\[
  A_n^{(M)}=\{X_n\ge M n\log_2 n\}.
\]
Then
\[
  \Prob(A_n^{(M)})\ge \frac{c_M}{n\log_2 n}
\]
for large $n$, and the series over $n$ diverges.  The events
$A_n^{(M)}$ are independent, so by the second Borel-Cantelli lemma they
occur infinitely often a.s.  Hence
\[
  \limsup_{n\to\infty}\frac{X_n}{n\log_2 n}\ge M
  \qquad\text{a.s.}
\]
Intersecting over integer $M\ge1$ gives the limsup equal to infinity.
Since $S_n\ge X_n$, the same conclusion holds for
$S_n/(n\log_2 n)$.

\section*{Exercise 2.4.1}

\noindent\textbf{Question.}
Lazy janitor.  Suppose the $i$th light bulb burns for an amount of time
$X_i$ and then remains burned out for time $Y_i$ before being replaced.
Suppose the $X_i,Y_i$ are positive and independent with the $X$'s having
distribution $F$ and the $Y$'s having distribution $G$, both of which
have finite mean.  Let $R_t$ be the amount of time in $[0,t]$ that we
have a working light bulb.  Show that
\[
  R_t/t\to \frac{\E X_i}{\E X_i+\E Y_i}
  \qquad\text{a.s.}
\]

\Knowledge{
\path{durrett_ch2_iid_finite_mean_slln};
\path{durrett_ch2_renewal_process};
\path{durrett_ch2_slln_normalized_sums}.
}

\noindent\textbf{Solution.}
Let $Z_i=X_i+Y_i$ be the length of the $i$th complete cycle and
$T_n=Z_1+\cdots+Z_n$.  Let
\[
  N(t)=\max\{n:T_n\le t\}.
\]
By the strong law,
\[
  \frac{T_n}{n}\to \E Z_1=\E X_1+\E Y_1,
  \qquad
  \frac{X_1+\cdots+X_n}{n}\to \E X_1
  \qquad\text{a.s.}
\]
Thus $N(t)\to\infty$ and
\[
  \frac{N(t)}{t}\to \frac{1}{\E X_1+\E Y_1}
  \qquad\text{a.s.}
\]
The working time up to $t$ differs from
$X_1+\cdots+X_{N(t)}$ by at most the working part of one incomplete
cycle.  Hence
\[
  \sum_{i=1}^{N(t)}X_i
  \le R_t
  \le \sum_{i=1}^{N(t)+1}X_i.
\]
Dividing by $t$ and using the two strong-law limits gives
\[
  \frac{R_t}{t}
  \to
  \E X_1\cdot \frac{1}{\E X_1+\E Y_1}
  =
  \frac{\E X_1}{\E X_1+\E Y_1}
  \qquad\text{a.s.}
\]

\section*{Exercise 2.4.2}

\noindent\textbf{Question.}
Let $X_0=(1,0)$ and define $X_n\in\R^2$ inductively by declaring that
$X_{n+1}$ is chosen at random from the ball of radius $|X_n|$ centered
at the origin, i.e. $X_{n+1}/|X_n|$ is uniformly distributed on the ball
of radius $1$ and independent of $X_1,\ldots,X_n$.  Prove that
\[
  n^{-1}\log |X_n|\to c
\]
almost surely and compute $c$.

\Knowledge{
\path{durrett_ch2_iid_finite_mean_slln};
\path{durrett_ch2_functions_of_independent_blocks};
\path{durrett_ch2_slln_normalized_sums}.
}

\noindent\textbf{Solution.}
Write
\[
  X_{n+1}=|X_n|U_{n+1},
\]
where $U_1,U_2,\ldots$ are i.i.d. uniform random vectors in the unit disk.
Then
\[
  |X_n|=\prod_{k=1}^n |U_k|,
  \qquad
  \log|X_n|=\sum_{k=1}^n \log|U_k|.
\]
If $R=|U_1|$, then $R$ has density $2r$ on $0\le r\le1$.  Hence
\[
  \E|\log R|
  =\int_0^1 2r|\log r|\,dr<\infty,
\]
and the strong law applies:
\[
  \frac{1}{n}\log|X_n|
  \to \E\log R
  =\int_0^1 2r\log r\,dr
  =-\frac12.
\]
Thus $c=-1/2$.

\section*{Exercise 2.4.3}

\noindent\textbf{Question.}
Investment problem.  At the beginning of each year you can buy bonds
for 1 unit that are worth $a$ at the end of the year or stocks that are
worth a random amount $V\ge0$.  If you always invest a fixed proportion
$p$ of your wealth in bonds, then your wealth at the end of year
$n+1$ is
\[
  W_{n+1}=(ap+(1-p)V_n)W_n.
\]
Suppose $V_1,V_2,\ldots$ are i.i.d. with $\E V_n^2<\infty$ and
$\E(V_n^{-2})<\infty$.
\begin{enumerate}[label=(\roman*)]
\item Show that $n^{-1}\log W_n\to c(p)$ a.s.
\item Show that $c(p)$ is concave.
\item By investigating $c'(0)$ and $c'(1)$, give conditions on $V$ that
guarantee that the optimal choice of $p$ is in $(0,1)$.
\item Suppose $\Prob(V=1)=\Prob(V=4)=1/2$.  Find the optimal $p$ as a
function of $a$.
\end{enumerate}

\Knowledge{
\path{durrett_ch2_iid_finite_mean_slln};
\path{durrett_ch1_jensen_integral};
\path{durrett_ch2_slln_normalized_sums}.
}

\noindent\textbf{Solution.}
Assume $0\le p\le1$ and $a>0$.  Iterating the wealth recursion,
\[
  \log W_n
  =\log W_0+\sum_{k=0}^{n-1}\log(ap+(1-p)V_k).
\]
The assumptions $\E V^2<\infty$ and $\E V^{-2}<\infty$ imply that the
positive and negative parts of $\log(ap+(1-p)V)$ are integrable for
$0\le p\le1$.  Therefore the strong law gives
\[
  \frac{1}{n}\log W_n\to
  c(p):=\E\log(ap+(1-p)V)
  \qquad\text{a.s.}
\]

For each fixed $v>0$, the function
\[
  p\mapsto \log(ap+(1-p)v)
\]
is concave because it is the logarithm of an affine positive function.
Taking expectations preserves concavity, so $c(p)$ is concave.

Differentiating under the expectation,
\[
  c'(p)=\E\left[\frac{a-V}{ap+(1-p)V}\right].
\]
Thus
\[
  c'(0)=a\,\E(1/V)-1,
  \qquad
  c'(1)=1-\frac{\E V}{a}.
\]
Since $c$ is concave, an interior maximizer is guaranteed when
\[
  c'(0)>0
  \quad\text{and}\quad
  c'(1)<0,
\]
that is,
\[
  a\,\E(1/V)>1
  \quad\text{and}\quad
  \E V>a.
\]

Now suppose $\Prob(V=1)=\Prob(V=4)=1/2$.  Then
\[
  c(p)=\frac12\log(1+(a-1)p)
       +\frac12\log(4+(a-4)p).
\]
The endpoint derivative conditions are
\[
  c'(0)=\frac{5a}{8}-1,\qquad
  c'(1)=1-\frac{5}{2a}.
\]
Hence the optimum is $p=0$ if $a\le8/5$, and $p=1$ if
$a\ge5/2$.  For $8/5<a<5/2$, solve $c'(p)=0$:
\[
  (a-1)(4+(a-4)p)+(a-4)(1+(a-1)p)=0.
\]
Thus
\[
  p^*(a)=\frac{8-5a}{2(a-1)(a-4)},
  \qquad 8/5<a<5/2.
\]
So
\[
  p^*(a)=
  \begin{cases}
  0, & a\le 8/5,\\[3pt]
  \dfrac{8-5a}{2(a-1)(a-4)}, & 8/5<a<5/2,\\[8pt]
  1, & a\ge 5/2.
  \end{cases}
\]

\section*{Exercise 2.5.1}

\noindent\textbf{Question.}
Suppose $X_i$ are i.i.d., $\E X_i=0$, and $\operatorname{var}(X_i)=C<\infty$.
Use Kolmogorov's maximal inequality with $n=m^\alpha$, where
$\alpha(2p-1)>1$, to show that if $S_n=\sum_{i=1}^n X_i$ and $p>1/2$,
then
\[
  \frac{S_n}{n^p}\to0\qquad\text{a.s.}
\]

\Knowledge{
\path{durrett_ch2_kolmogorov_maximal_inequality};
\path{durrett_ch2_borel_cantelli};
\path{durrett_ch2_subsequence_upgrade}.
}

\noindent\textbf{Solution.}
Let $n_m=\lfloor m^\alpha\rfloor$, with $\alpha(2p-1)>1$.  By
Kolmogorov's maximal inequality,
\[
  \Prob\left(\max_{k\le n_m}|S_k|>\varepsilon n_m^p\right)
  \le {C n_m\over \varepsilon^2 n_m^{2p}}
  ={C\over\varepsilon^2}n_m^{1-2p}.
\]
Since $n_m^{1-2p}\le c m^{-\alpha(2p-1)}$ and
$\alpha(2p-1)>1$, these probabilities are summable.  Borel--Cantelli
therefore gives
\[
  \max_{k\le n_m}{|S_k|\over n_m^p}\to0 \qquad\text{a.s.}
\]
If $n_m<k\le n_{m+1}$, then
\[
  {|S_k|\over k^p}\le
  \left({n_{m+1}\over n_m}\right)^p
  \max_{j\le n_{m+1}}{|S_j|\over n_{m+1}^p}.
\]
Because $n_{m+1}/n_m\to1$, the right side tends to $0$ a.s.  Hence
$S_n/n^p\to0$ a.s.

\section*{Exercise 2.5.2}

\noindent\textbf{Question.}
Let $p>0$.  Prove the converse of Theorem 2.5.12: if
$S_n/n^{1/p}\to0$ a.s. for i.i.d. $X_i$, then
$\E|X_1|^p<\infty$.

\Knowledge{
\path{durrett_ch2_borel_cantelli};
\path{durrett_ch2_tail_sum_moment_test};
\path{durrett_ch2_slln_normalized_sums}.
}

\noindent\textbf{Solution.}
Since
\[
  {X_n\over n^{1/p}}
  ={S_n\over n^{1/p}}
   -\left({n-1\over n}\right)^{1/p}
    {S_{n-1}\over (n-1)^{1/p}},
\]
we have $X_n/n^{1/p}\to0$ a.s.  Hence for every $\varepsilon>0$,
\[
  \Prob(|X_n|>\varepsilon n^{1/p}\ \text{i.o.})=0.
\]
The events are independent, so the second Borel--Cantelli lemma implies
\[
  \sum_{n=1}^\infty
  \Prob(|X_1|^p>\varepsilon^p n)<\infty.
\]
By the tail-sum criterion this is equivalent to
$\E|X_1|^p<\infty$.

\section*{Exercise 2.5.3}

\noindent\textbf{Question.}
Let $X_i$ be independent standard normal random variables.  Show that,
for every fixed $t$,
\[
  \sum_{n=1}^\infty {X_n\sin(n\pi t)\over n}
\]
converges a.s.

\Knowledge{
\path{durrett_ch2_two_series_theorem};
\path{durrett_ch2_independent_series_variance_test}.
}

\noindent\textbf{Solution.}
For fixed $t$, set
\[
  Y_n={X_n\sin(n\pi t)\over n}.
\]
Then the $Y_n$ are independent, $\E Y_n=0$, and
\[
  \sum_{n=1}^\infty \operatorname{var}(Y_n)
  =\sum_{n=1}^\infty {\sin^2(n\pi t)\over n^2}
  \le \sum_{n=1}^\infty {1\over n^2}<\infty.
\]
The independent-series variance test therefore gives
$\sum_n Y_n$ converges a.s.

\section*{Exercise 2.5.4}

\noindent\textbf{Question.}
Let $X_n$ be independent with $\E X_n=0$ and
$\operatorname{var}(X_n)=\sigma_n^2$.
\begin{enumerate}[label=(\roman*)]
\item If $\sum_n \sigma_n^2/n^2<\infty$, show that
$\sum_n X_n/n$ converges a.s., and hence
$n^{-1}\sum_{m=1}^n X_m\to0$ a.s.
\item If $\sum_n \sigma_n^2/n^2=\infty$ and $\sigma_n^2\le n^2$,
construct independent $X_n$ with $\E X_n=0$ and
$\operatorname{var}(X_n)\le\sigma_n^2$ such that $X_n/n$, and hence
$n^{-1}\sum_{m=1}^nX_m$, does not converge to $0$ a.s.
\end{enumerate}

\Knowledge{
\path{durrett_ch2_two_series_theorem};
\path{durrett_ch2_kronecker_lemma};
\path{durrett_ch2_borel_cantelli}.
}

\noindent\textbf{Solution.}
(i) The variables $X_n/n$ are independent, have mean $0$, and
\[
  \sum_n \operatorname{var}(X_n/n)
  =\sum_n {\sigma_n^2\over n^2}<\infty.
\]
Thus $\sum_n X_n/n$ converges a.s.  Kronecker's lemma gives
\[
  {1\over n}\sum_{m=1}^n X_m\to0\qquad\text{a.s.}
\]

(ii) Let
\[
  \Prob(X_n=n)=\Prob(X_n=-n)={\sigma_n^2\over 2n^2},
  \qquad
  \Prob(X_n=0)=1-{\sigma_n^2\over n^2}.
\]
Then $\E X_n=0$ and $\operatorname{var}(X_n)=\sigma_n^2$.  Also
\[
  \sum_n \Prob(|X_n/n|=1)=\sum_n {\sigma_n^2\over n^2}=\infty.
\]
By Borel--Cantelli, $|X_n/n|=1$ infinitely often a.s.; hence
$X_n/n\not\to0$.  If the averages $S_n/n$ converged to $0$, then
\[
  {X_n\over n}={S_n\over n}-{n-1\over n}{S_{n-1}\over n-1}
\]
would also converge to $0$, a contradiction.

\section*{Exercise 2.5.5}

\noindent\textbf{Question.}
Let $X_n\ge0$ be independent.  Show that the following are equivalent:
\begin{enumerate}[label=(\roman*)]
\item $\sum_n X_n<\infty$ a.s.
\item $\sum_n\{\Prob(X_n>1)+\E[X_n\1_{\{X_n\le1\}}]\}<\infty$.
\item $\sum_n \E[X_n/(1+X_n)]<\infty$.
\end{enumerate}

\Knowledge{
\path{durrett_ch2_three_series_theorem};
\path{durrett_ch2_borel_cantelli};
\path{durrett_ch2_truncation_method}.
}

\noindent\textbf{Solution.}
For $x\ge0$,
\[
  {1\over2}\bigl(x\1_{\{x\le1\}}+\1_{\{x>1\}}\bigr)
  \le {x\over1+x}
  \le x\1_{\{x\le1\}}+\1_{\{x>1\}}.
\]
Thus (ii) and (iii) are equivalent.

If (ii) holds, then $\sum_n\Prob(X_n>1)<\infty$, so only finitely many
large jumps occur a.s.  Also
\[
  \E\sum_n X_n\1_{\{X_n\le1\}}
  =\sum_n\E[X_n\1_{\{X_n\le1\}}]<\infty,
\]
so the truncated sum is finite a.s.  Hence (i) holds.

Conversely assume (i), and put $Y_n=X_n/(1+X_n)$.  Then
$0\le Y_n\le1$ and $\sum_nY_n<\infty$ a.s.  If
$\sum_n\E Y_n=\infty$, then, using $1-e^{-y}\ge(1-e^{-1})y$ on
$0\le y\le1$,
\[
  \E e^{-\sum_n Y_n}
  =\prod_n \E e^{-Y_n}
  \le \prod_n \exp\{-(1-e^{-1})\E Y_n\}=0,
\]
which would force $\sum_nY_n=\infty$ a.s., a contradiction.  Therefore
$\sum_n\E Y_n<\infty$, which is (iii).

\section*{Exercise 2.5.6}

\noindent\textbf{Question.}
Let $\psi(x)=x^2$ for $|x|\le1$ and $\psi(x)=|x|$ for $|x|\ge1$.
Show that if $X_n$ are independent, $\E X_n=0$, and
$\sum_n\E\psi(X_n)<\infty$, then $\sum_nX_n$ converges a.s.

\Knowledge{
\path{durrett_ch2_three_series_theorem};
\path{durrett_ch2_truncation_method}.
}

\noindent\textbf{Solution.}
Use the three-series theorem with truncation at $1$.  First,
\[
  \sum_n\Prob(|X_n|>1)
  \le \sum_n\E[|X_n|\1_{\{|X_n|>1\}}]
  \le \sum_n\E\psi(X_n)<\infty.
\]
For $Y_n=X_n\1_{\{|X_n|\le1\}}$,
\[
  \sum_n\operatorname{var}(Y_n)
  \le \sum_n\E Y_n^2
  =\sum_n\E[X_n^2\1_{\{|X_n|\le1\}}]<\infty.
\]
Finally, since $\E X_n=0$,
\[
  |\E Y_n|
  =|\E[X_n\1_{\{|X_n|>1\}}]|
  \le \E[|X_n|\1_{\{|X_n|>1\}}],
\]
and the right side is summable.  All three conditions hold, so
$\sum_nX_n$ converges a.s.

\section*{Exercise 2.5.7}

\noindent\textbf{Question.}
Let $X_n$ be independent.  Suppose
\[
  \sum_n \E |X_n|^{p(n)}<\infty,
  \qquad 0<p(n)\le2,
\]
and $\E X_n=0$ when $p(n)>1$.  Show that $\sum_nX_n$ converges a.s.

\Knowledge{
\path{durrett_ch2_three_series_theorem};
\path{durrett_ch2_truncation_method}.
}

\noindent\textbf{Solution.}
For all $n$,
\[
  \Prob(|X_n|>1)\le \E|X_n|^{p(n)},
\]
so large jumps occur only finitely often a.s.  If $p(n)\le1$, then on
$\{|X_n|\le1\}$ we have $|X_n|\le |X_n|^{p(n)}$, so the corresponding
small-jump parts are absolutely summable in expectation, hence a.s.

For the indices with $p(n)>1$, apply the three-series theorem.  On
$|x|\le1$, $x^2\le |x|^{p(n)}$ because $p(n)\le2$, and on $|x|>1$,
$|x|\le |x|^{p(n)}$.  Hence
\[
  \sum_{p(n)>1}\operatorname{var}(X_n\1_{\{|X_n|\le1\}})<\infty
\]
and, using $\E X_n=0$,
\[
  \sum_{p(n)>1}
  |\E[X_n\1_{\{|X_n|\le1\}}]|
  \le \sum_{p(n)>1}\E[|X_n|\1_{\{|X_n|>1\}}]<\infty.
\]
Thus the full series converges a.s.

\section*{Exercise 2.5.8}

\noindent\textbf{Question.}
Let $X_n$ be i.i.d. and not identically $0$.  For the power series
$\sum_{n\ge1}X_n(\omega)z^n$, show that its radius of convergence is
$1$ a.s. if $\E\log^+|X_1|<\infty$, and is $0$ a.s. if
$\E\log^+|X_1|=\infty$.

\Knowledge{
\path{durrett_ch2_borel_cantelli};
\path{durrett_ch2_tail_sum_moment_test};
\path{durrett_ch2_power_series_radius}.
}

\noindent\textbf{Solution.}
The radius is
\[
  r(\omega)=\frac{1}{\limsup_n |X_n(\omega)|^{1/n}}.
\]
If $\E\log^+|X_1|<\infty$, then for every $\varepsilon>0$,
\[
  \sum_n \Prob(\log^+|X_n|>\varepsilon n)<\infty.
\]
Thus eventually $|X_n|^{1/n}\le e^\varepsilon$, and the limsup is at
most $1$.  Since $X_1$ is not identically $0$, there is $\delta>0$ with
$\Prob(|X_1|>\delta)>0$; by Borel--Cantelli this occurs infinitely
often, so the limsup is at least $1$.  Hence $r=1$.

If $\E\log^+|X_1|=\infty$, then for every $M>1$,
\[
  \sum_n\Prob(\log^+|X_n|>n\log M)=\infty.
\]
By Borel--Cantelli, $|X_n|>M^n$ infinitely often, so
$\limsup_n |X_n|^{1/n}\ge M$.  Since $M$ is arbitrary, the limsup is
infinite and $r=0$.

\section*{Exercise 2.5.9}

\noindent\textbf{Question.}
Let $X_n$ be independent and
$S_{m,n}=X_{m+1}+\cdots+X_n$.  Prove
\[
  \Prob\left(\max_{m<j\le n}|S_{m,j}|>2a\right)
  \min_{m<k\le n}\Prob(|S_{k,n}|\le a)
  \le \Prob(|S_{m,n}|>a).
\]

\Knowledge{
\path{durrett_ch2_levy_maximal_inequality};
\path{durrett_ch2_independent_increments}.
}

\noindent\textbf{Solution.}
Let $A_j$ be the event that $j$ is the first index in $(m,n]$ for which
$|S_{m,j}|>2a$.  The events $A_j$ are disjoint and
$A_j$ depends only on $X_{m+1},\ldots,X_j$, while $S_{j,n}$ depends only
on later variables.  If $A_j$ and $|S_{j,n}|\le a$ occur, then
\[
  |S_{m,n}|=|S_{m,j}+S_{j,n}|>a.
\]
Therefore
\[
  \Prob(A_j)\Prob(|S_{j,n}|\le a)
  =\Prob(A_j\cap\{|S_{j,n}|\le a\})
  \le \Prob(A_j\cap\{|S_{m,n}|>a\}).
\]
Summing over $j$ and taking the minimum over $k$ proves the claim.

\section*{Exercise 2.5.10}

\noindent\textbf{Question.}
Use the preceding inequality to prove Levy's theorem: if
$S_n=\sum_{i=1}^nX_i$ with independent $X_i$ and $\lim_nS_n$ exists in
probability, then $\lim_nS_n$ exists a.s.

\Knowledge{
\path{durrett_ch2_levy_maximal_inequality};
\path{durrett_ch2_cauchy_in_probability};
\path{durrett_ch2_borel_cantelli}.
}

\noindent\textbf{Solution.}
Since $S_n$ converges in probability, it is Cauchy in probability:
for every $a>0$,
\[
  \sup_{k,n\ge m}\Prob(|S_n-S_k|>a)\to0.
\]
Applying Exercise 2.5.9 to the block $(m,n]$, for large $m$ the
minimum factor is at least $1/2$, hence
\[
  \Prob\left(\max_{m<j\le n}|S_j-S_m|>2a\right)
  \le 2\Prob(|S_n-S_m|>a).
\]
Thus the partial sums are Cauchy in probability uniformly over each
tail maximum.  Choose $m_r\uparrow\infty$ so that
\[
  \Prob\left(\max_{m_r<j\le m_{r+1}}
  |S_j-S_{m_r}|>2^{-r}\right)<2^{-r}.
\]
Borel--Cantelli implies that, a.s., only finitely many of these block
oscillations exceed $2^{-r}$.  Since $S_{m_r}$ is Cauchy a.s. along the
chosen subsequence, the full sequence $S_n$ is Cauchy a.s. and therefore
converges a.s.

\section*{Exercise 2.5.11}

\noindent\textbf{Question.}
Let $X_i$ be i.i.d. and $S_n=\sum_{i=1}^nX_i$.  Use the inequality in
Exercise 2.5.9 to show that if $S_n/n\to0$ in probability, then
\[
  {1\over n}\max_{1\le m\le n}S_m\to0
  \qquad\text{in probability.}
\]

\Knowledge{
\path{durrett_ch2_levy_maximal_inequality};
\path{durrett_ch2_convergence_in_probability}.
}

\noindent\textbf{Solution.}
It is enough to prove the stronger statement with $\max |S_m|$.  Put
$m=0$ in Exercise 2.5.9.  For fixed $\varepsilon>0$,
\[
  \Prob\left(\max_{1\le j\le n}|S_j|>2\varepsilon n\right)
  \min_{1\le k\le n}\Prob(|S_n-S_k|\le \varepsilon n)
  \le \Prob(|S_n|>\varepsilon n).
\]
Since $S_n/n\to0$ in probability, the right side tends to $0$.  Also
$S_n-S_k$ has the same law as $S_{n-k}$, and
\[
  \min_{0\le \ell\le n}\Prob(|S_\ell|\le \varepsilon n)\to1;
\]
for large $\ell$ this follows from $|S_\ell|/\ell\to0$ in probability,
and for finitely many small $\ell$ it is immediate because
$\varepsilon n\to\infty$.  Hence
\[
  n^{-1}\max_{1\le j\le n}|S_j|\to0
  \quad\text{in probability},
\]
which implies the desired one-sided result.

\section*{Exercise 2.5.12}

\noindent\textbf{Question.}
Let $X_i$ be i.i.d., $S_n=\sum_{i=1}^nX_i$, and let $a(n)\uparrow\infty$
with $a(2^n)/a(2^{n-1})$ bounded.
\begin{enumerate}[label=(\roman*)]
\item Use Exercise 2.5.9 to show that if $S_n/a(n)\to0$ in probability
and $S_{2^n}/a(2^n)\to0$ a.s., then $S_n/a(n)\to0$ a.s.
\item If $\E X_1=0$ and $\E X_1^2<\infty$, conclude that for every
$\varepsilon>0$,
\[
  {S_n\over n^{1/2}(\log_2 n)^{1/2+\varepsilon}}\to0
  \qquad\text{a.s.}
\]
\end{enumerate}

\Knowledge{
\path{durrett_ch2_levy_maximal_inequality};
\path{durrett_ch2_chebyshev};
\path{durrett_ch2_borel_cantelli}.
}

\noindent\textbf{Solution.}
(i) Write $b_k=a(2^k)$.  On the dyadic block
$2^{k-1}<j\le2^k$, Exercise 2.5.9 gives, for fixed $\eta>0$,
\[
  \Prob\left(\max_{2^{k-1}<j\le2^k}|S_j-S_{2^{k-1}}|>2\eta b_k\right)
  \le
  { \Prob(|S_{2^k}-S_{2^{k-1}}|>\eta b_k)\over q_k},
\]
where
\[
  q_k=\min_{2^{k-1}<j\le2^k}
  \Prob(|S_{2^k}-S_j|\le \eta b_k).
\]
The assumption $S_n/a(n)\to0$ in probability and the boundedness of
$a(2^k)/a(2^{k-1})$ imply $q_k\to1$.  The dyadic convergence
$S_{2^k}/b_k\to0$ a.s. makes the endpoint increments negligible; the
standard Borel--Cantelli blocking argument from Levy's theorem then
gives
\[
  \max_{2^{k-1}<j\le2^k}{|S_j-S_{2^{k-1}}|\over b_k}\to0
  \qquad\text{a.s.}
\]
Since $a(j)$ is increasing and $b_k/a(j)$ is bounded on the block, and
since $S_{2^{k-1}}/a(2^{k-1})\to0$ a.s., it follows that
$S_n/a(n)\to0$ a.s.

(ii) Let
\[
  a(n)=n^{1/2}(\log_2 n)^{1/2+\varepsilon},\qquad n\ge2.
\]
Then $a(2^k)=2^{k/2}k^{1/2+\varepsilon}$ and
\[
  \sum_{k=1}^\infty
  \Prob(|S_{2^k}|>\delta a(2^k))
  \le
  \sum_{k=1}^\infty
  {\operatorname{var}(S_{2^k})\over \delta^2 a(2^k)^2}
  =
  { \operatorname{var}(X_1)\over\delta^2}
  \sum_{k=1}^\infty {1\over k^{1+2\varepsilon}}<\infty.
\]
Borel--Cantelli yields $S_{2^k}/a(2^k)\to0$ a.s.  Chebyshev also gives
$S_n/a(n)\to0$ in probability.  Part (i) completes the proof.

\section*{Exercise 2.6.1}

\noindent\textbf{Question.}
For a renewal process with interarrival times $\xi_i$, renewal count
$N_t$, and renewal function $U(t)=\E N_t$, prove
\[
  {t\over \E(\xi_1\wedge t)}
  \le U(t)\le
  {2t\over \E(\xi_1\wedge t)}.
\]

\Knowledge{
\path{durrett_ch2_renewal_function};
\path{durrett_ch2_wald_identity_truncated}.
}

\noindent\textbf{Solution.}
Let $\eta_i=\xi_i\wedge t$.  Since $T_{N_t}>t$, while each
$\eta_i\le t$,
\[
  t\le \sum_{i=1}^{N_t}\eta_i\le 2t.
\]
Indeed, the partial sum before the last term is at most $t$, and the
last truncated term is at most $t$.  Wald's identity for the bounded
increments $\eta_i$ gives
\[
  \E\sum_{i=1}^{N_t}\eta_i=\E(\xi_1\wedge t)\,\E N_t.
\]
Dividing by $\E(\xi_1\wedge t)$ proves the two inequalities.

\section*{Exercise 2.6.2}

\noindent\textbf{Question.}
Deduce the elementary renewal theorem from the strong law by showing
\[
  \limsup_{t\to\infty}\E (N_t/t)^2<\infty.
\]

\Knowledge{
\path{durrett_ch2_elementary_renewal_theorem};
\path{durrett_ch2_renewal_slln};
\path{durrett_ch2_uniform_integrability}.
}

\noindent\textbf{Solution.}
Choose $\delta>0$ with $p=\Prob(\xi_1>\delta)>0$, and let
$K=\lceil t/\delta\rceil$.  If $N_t>mK$, then among the first $mK$
interarrival times no block of length $K$ can consist entirely of
variables larger than $\delta$.  Thus
\[
  \Prob(N_t>mK)\le (1-p^K)^m.
\]
This geometric tail bound implies $\E N_t^2\le C K^2\le C't^2$, so
$\limsup_t\E(N_t/t)^2<\infty$.

By the strong law for renewal times, $N_t/t\to1/\mu$ a.s. when
$\mu=\E\xi_1<\infty$.  The $L^2$ bound gives uniform integrability, so
$\E(N_t/t)\to1/\mu$, i.e. $U(t)/t\to1/\mu$.

\section*{Exercise 2.6.3}

\noindent\textbf{Question.}
Poisson arrivals occur at rate $1$.  A server accepts a customer iff the
server is free, and service times have distribution $F$ with mean $\mu$.
Show that accepted arrival times form a renewal process with mean
interarrival time $\mu+1$, and that the limiting fraction of arrivals
served is $1/(\mu+1)$.

\Knowledge{
\path{durrett_ch2_poisson_memoryless};
\path{durrett_ch2_renewal_reward_theorem};
\path{durrett_ch2_renewal_slln}.
}

\noindent\textbf{Solution.}
After an accepted customer begins service, the server is unavailable for
a service time with mean $\mu$.  Once service ends, the memoryless
property of the Poisson process says the waiting time to the next
arrival is exponential with mean $1$, independent of the past.  Hence
accepted arrivals form a renewal process with interarrival law
\[
  \text{service time}+\operatorname{Exp}(1),
\]
whose mean is $\mu+1$.  Therefore the accepted arrival rate is
$1/(\mu+1)$.  The total arrival rate is $1$, so the limiting fraction
served is $1/(\mu+1)$.

\section*{Exercise 2.6.4}

\noindent\textbf{Question.}
Let $A_t=t-T_{N(t)-1}$ be the age.  For fixed $x>0$, show that
$H(t)=\Prob(A_t>x)$ satisfies
\[
  H(t)=(1-F(t))\1_{(x,\infty)}(t)+\int_0^t H(t-s)\,dF(s),
\]
and conclude
\[
  \Prob(A_t>x)\to {1\over\mu}\int_x^\infty (1-F(u))\,du.
\]

\Knowledge{
\path{durrett_ch2_renewal_equation};
\path{durrett_ch2_key_renewal_theorem};
\path{durrett_ch2_age_residual_life}.
}

\noindent\textbf{Solution.}
Condition on the first renewal time $T_1$.  If $T_1>t$, then
$A_t=t$, so the event occurs exactly when $t>x$.  If $T_1=s\le t$, the
process restarts at time $s$, and the new age at time $t$ has the same
law as $A_{t-s}$.  This gives the stated renewal equation.

The forcing function is
\[
  h(t)=(1-F(t))\1_{(x,\infty)}(t),
\]
which is directly Riemann integrable when $\mu<\infty$.  The key renewal
theorem yields
\[
  H(t)\to {1\over\mu}\int_0^\infty h(u)\,du
  ={1\over\mu}\int_x^\infty(1-F(u))\,du.
\]

\section*{Exercise 2.6.5}

\noindent\textbf{Question.}
Use the renewal equation in Exercise 2.6.4 and Theorem 2.6.9 to show
that if the renewal process is a rate $\lambda$ Poisson process, then
$A_t$ has the same distribution as $\xi_1\wedge t$.

\Knowledge{
\path{durrett_ch2_poisson_memoryless};
\path{durrett_ch2_renewal_equation};
}

\noindent\textbf{Solution.}
For Poisson renewals, $F(u)=1-e^{-\lambda u}$.  For $0<x<t$,
\[
  \Prob(A_t>x)=\Prob(\text{no renewal in }(t-x,t])=e^{-\lambda x},
\]
and for $x\ge t$ this probability is $0$.  This is also
\[
  \Prob(\xi_1\wedge t>x),
\]
since $\xi_1$ is exponential with rate $\lambda$.  Thus
$A_t\stackrel d=\xi_1\wedge t$.

\section*{Exercise 2.6.6}

\noindent\textbf{Question.}
Let $B_t=T_{N(t)}-t$ be the residual lifetime.  Show that
\[
  \Prob(A_t>x,\ B_t>y)
  \to {1\over\mu}\int_{x+y}^\infty(1-F(u))\,du.
\]

\Knowledge{
\path{durrett_ch2_age_residual_life};
\path{durrett_ch2_key_renewal_theorem};
}

\noindent\textbf{Solution.}
The event $\{A_t>x,B_t>y\}$ means that the renewal interval straddling
$t$ began before $t-x$ and has length exceeding $x+y$.  Its renewal
equation has forcing term
\[
  h(t)=(1-F(t+y))\1_{(x,\infty)}(t).
\]
By the key renewal theorem,
\[
  \Prob(A_t>x,B_t>y)
  \to {1\over\mu}\int_0^\infty h(u)\,du
  ={1\over\mu}\int_x^\infty(1-F(u+y))\,du
  ={1\over\mu}\int_{x+y}^\infty(1-F(v))\,dv.
\]

\section*{Exercise 2.6.7}

\noindent\textbf{Question.}
In an alternating renewal process, working times have distribution
$F_1$ with mean $\mu_1$, repair times have distribution $F_2$ with mean
$\mu_2$, and one cycle has law $F=F_1*F_2$.  If $F$ is nonarithmetic,
show that the probability $H(t)$ that the machine is working at time
$t$ satisfies
\[
  H(t)\to {\mu_1\over\mu_1+\mu_2}.
\]

\Knowledge{
\path{durrett_ch2_alternating_renewal_process};
\path{durrett_ch2_key_renewal_theorem};
\path{durrett_ch2_renewal_reward_theorem}.
}

\noindent\textbf{Solution.}
The machine is working at time $t$ either because the initial working
period exceeds $t$, or because a full cycle has ended and the restarted
process is working.  Thus
\[
  H(t)=1-F_1(t)+\int_0^t H(t-s)\,dF(s).
\]
The forcing function $h(t)=1-F_1(t)$ is directly Riemann integrable and
\[
  \int_0^\infty h(t)\,dt=\mu_1.
\]
The mean cycle length is $\mu_1+\mu_2$.  The key renewal theorem gives
\[
  H(t)\to {\mu_1\over \mu_1+\mu_2}.
\]

\section*{Exercise 2.6.8}

\noindent\textbf{Question.}
Find a renewal equation for
$H(t)=\Prob(\text{the number of renewals in }[0,t]\text{ is odd})$, and
use the renewal theorem to show that $H(t)\to1/2$.

\Knowledge{
\path{durrett_ch2_renewal_equation};
\path{durrett_ch2_alternating_renewal_process};
\path{durrett_ch2_key_renewal_theorem}.
}

\noindent\textbf{Solution.}
Counting the renewal at time $0$, if the first interarrival exceeds
$t$, the count is $1$ and hence odd.  If the first interarrival is
$s\le t$, parity is reversed relative to the restarted process.  Hence
\[
  H(t)=1-F(t)+\int_0^t(1-H(t-s))\,dF(s).
\]
This is the alternating-renewal equation with equal on- and off-cycle
distributions.  By Exercise 2.6.7, the limiting fraction of time spent
in either parity is
\[
  {\mu\over\mu+\mu}={1\over2}.
\]
Thus $H(t)\to1/2$.

\section*{Exercise 2.6.9}

\noindent\textbf{Question.}
Assume $F$ has directly Riemann integrable density $f$.  Let
$V=U-\1_{[0,\infty)}$.  Show that $V$ has density $v$ satisfying
\[
  v(t)=f(t)+\int_0^t v(t-s)\,dF(s),
\]
and conclude that $v(t)\to1/\mu$.

\Knowledge{
\path{durrett_ch2_renewal_density};
\path{durrett_ch2_key_renewal_theorem}.
}

\noindent\textbf{Solution.}
Since
\[
  U=\sum_{n=0}^\infty F^{*n},
  \qquad
  V=\sum_{n=1}^\infty F^{*n},
\]
the measure $V$ has density
\[
  v(t)=\sum_{n=1}^\infty f^{*n}(t).
\]
Separating the first term gives
\[
  v(t)=f(t)+\sum_{n=2}^\infty f^{*n}(t)
  =f(t)+\int_0^t v(t-s)\,dF(s).
\]
The key renewal theorem applied to the directly Riemann integrable
function $f$ gives
\[
  v(t)\to {1\over\mu}\int_0^\infty f(s)\,ds={1\over\mu}.
\]

\section*{Exercise 2.7.1}

\noindent\textbf{Question.}
For $\gamma(a)$ in the large-deviation theorem, show that the following
are equivalent:
\[
  \gamma(a)=-\infty,\qquad
  \Prob(X_1\ge a)=0,\qquad
  \Prob(S_n\ge na)=0\text{ for all }n.
\]

\Knowledge{
\path{durrett_ch2_large_deviation_gamma};
\path{durrett_ch2_independent_sums}.
}

\noindent\textbf{Solution.}
If $\Prob(X_1\ge a)=0$, then every summand is $<a$ a.s.; hence
$S_n<na$ a.s. and $\Prob(S_n\ge na)=0$ for all $n$, so
$\gamma(a)=-\infty$.

Conversely, if $\Prob(S_n\ge na)=0$ for all $n$, taking $n=1$ gives
$\Prob(X_1\ge a)=0$.  Finally, if $\gamma(a)=-\infty$ but
$p=\Prob(X_1\ge a)>0$, then
$\Prob(S_n\ge na)\ge p^n$, so
$n^{-1}\log\Prob(S_n\ge na)\ge\log p>-\infty$, a contradiction.

\section*{Exercise 2.7.2}

\noindent\textbf{Question.}
Use the definition of $\gamma$ to show that for rational
$\lambda\in[0,1]$,
\[
  \gamma(\lambda a+(1-\lambda)b)
  \ge \lambda\gamma(a)+(1-\lambda)\gamma(b).
\]
Extend this to all $\lambda$ and show that $\gamma$ is concave, hence
Lipschitz continuous on compact subsets of $\{\gamma>-\infty\}$.

\Knowledge{
\path{durrett_ch2_large_deviation_gamma};
\path{durrett_ch2_superadditivity};
\path{durrett_ch2_concavity}.
}

\noindent\textbf{Solution.}
Let $\lambda=r/(r+s)$ with integers $r,s\ge0$.  Independence gives
\[
  \Prob(S_{r+s}\ge ra+sb)
  \ge \Prob(S_r\ge ra)\Prob(S_s'\ge sb),
\]
where $S_s'$ is an independent copy of $S_s$.  Iterating this block
argument and taking logarithms divided by the total length yields
\[
  \gamma\!\left({ra+sb\over r+s}\right)
  \ge {r\over r+s}\gamma(a)+{s\over r+s}\gamma(b).
\]
Thus the inequality holds for rational $\lambda$.  Since $\gamma$ is
nonincreasing in $a$, approximating an arbitrary $\lambda$ from above
and below by rationals extends the inequality to all $\lambda\in[0,1]$.
Therefore $\gamma$ is concave on its effective domain.  A finite concave
function on an open interval has bounded one-sided slopes on compact
subintervals, hence is Lipschitz there.

\section*{Exercise 2.7.3}

\noindent\textbf{Question.}
Let $X_i$ be i.i.d. Poisson random variables with mean $1$.  Find
\[
  \lim_{n\to\infty}{1\over n}\log\Prob(S_n\ge na),
  \qquad a>1.
\]

\Knowledge{
\path{durrett_ch2_cramer_theorem};
\path{durrett_ch2_poisson_large_deviation};
\path{durrett_ch2_chernoff_bound}.
}

\noindent\textbf{Solution.}
Here $S_n$ is Poisson with mean $n$.  The moment generating function of
$X_1$ is
\[
  \phi(\theta)=\exp(e^\theta-1).
\]
For $a>1$, the optimizing Chernoff parameter solves
$\phi'(\theta)/\phi(\theta)=e^\theta=a$, so $\theta=\log a$.  The rate
function is
\[
  I(a)=a\log a-a+1.
\]
Therefore
\[
  \lim_{n\to\infty}{1\over n}\log\Prob(S_n\ge na)
  =-I(a)
  =-a\log a+a-1.
\]

\section*{Exercise 2.7.4}

\noindent\textbf{Question.}
For fair coin flips written as $X_i=\pm1$, prove
\[
  \phi(\theta)\le \exp(\phi(\theta)-1)\le \exp(\beta\theta^2),
  \qquad \theta\le1,
\]
where
\[
  \beta=\sum_{n=1}^\infty {1\over(2n)!}\approx0.586,
\]
and use Chernoff's inequality to show that for $0\le a\le1$,
\[
  \Prob(S_n\ge an)\le \exp\left(-{na^2\over4\beta}\right).
\]

\Knowledge{
\path{durrett_ch2_chernoff_bound};
\path{durrett_ch2_mgf_bounds};
\path{durrett_ch2_subgaussian_coin_flips}.
}

\noindent\textbf{Solution.}
For $X_1=\pm1$ with probability $1/2$,
\[
  \phi(\theta)=\E e^{\theta X_1}=\cosh\theta.
\]
Since $x\le e^{x-1}$ for $x>0$,
\[
  \phi(\theta)\le e^{\phi(\theta)-1}.
\]
For $0\le\theta\le1$,
\[
  \phi(\theta)-1
  =\sum_{n=1}^\infty {\theta^{2n}\over(2n)!}
  \le \theta^2\sum_{n=1}^\infty {1\over(2n)!}
  =\beta\theta^2.
\]
Thus $\phi(\theta)\le e^{\beta\theta^2}$.  Chernoff's bound gives
\[
  \Prob(S_n\ge an)
  \le \exp\{-n(a\theta-\beta\theta^2)\}.
\]
For $0\le a\le1$, choose $\theta=a/(2\beta)\le1$.  Then
$a\theta-\beta\theta^2=a^2/(4\beta)$, so
\[
  \Prob(S_n\ge an)\le
  \exp\left(-{na^2\over4\beta}\right).
\]

\section*{Exercise 2.7.5}

\noindent\textbf{Question.}
Suppose $\E X_i=0$ and $\E e^{\theta X_i}=\infty$ for all
$\theta>0$.  Show that for every $a>0$,
\[
  {1\over n}\log\Prob(S_n\ge na)\to0.
\]

\Knowledge{
\path{durrett_ch2_heavy_tail_large_deviation};
\path{durrett_ch2_one_big_jump};
\path{durrett_ch2_large_deviation_gamma}.
}

\noindent\textbf{Solution.}
The limit $\gamma(a)$ exists by the large-deviation theorem and is at
most $0$.  Suppose $\gamma(a)<0$ for some $a>0$.  Then for some
$c>0$,
\[
  \Prob(S_n\ge na)\le e^{-cn}
\]
for all large $n$.  By the weak law, for large $n$,
\[
  \Prob(S_{n-1}\ge -na)\ge {1\over2}.
\]
Therefore
\[
  \Prob(X_1\ge 2na)
  \le 2\Prob(S_n\ge na)
  \le 2e^{-cn}.
\]
Monotonicity fills the gaps between multiples of $2a$, so the positive
tail of $X_1$ is exponentially bounded.  This would imply
$\E e^{\theta X_1}<\infty$ for some $\theta>0$, contradicting the
assumption.  Hence $\gamma(a)=0$ for every $a>0$.

\section*{Exercise 2.7.6}

\noindent\textbf{Question.}
Show that if $\E X_i=0$, then for $a,\varepsilon>0$,
\[
  \liminf_{n\to\infty}
  {\Prob(S_n\ge na)\over n\Prob(X_1\ge n(a+\varepsilon))}
  \ge1.
\]

\Knowledge{
\path{durrett_ch2_one_big_jump};
\path{durrett_ch2_weak_law};
\path{durrett_ch2_tail_probability}.
}

\noindent\textbf{Solution.}
Let $p_n=\Prob(X_1\ge n(a+\varepsilon))$.  Since $\E|X_1|<\infty$,
$np_n\to0$.  Consider the event that exactly one of
$X_1,\ldots,X_n$ is at least $n(a+\varepsilon)$ and the remaining sum is
at least $-n\varepsilon$.  On this event, $S_n\ge na$.  Independence
and the weak law give
\[
  \Prob(S_{n-1}\ge -n\varepsilon)\to1.
\]
Hence
\[
  \Prob(S_n\ge na)
  \ge n p_n(1-p_n)^{n-1}\Prob(S_{n-1}\ge -n\varepsilon),
\]
and the right side is asymptotic to $n p_n$.  Dividing by
$n\Prob(X_1\ge n(a+\varepsilon))$ proves the desired lower bound.

\end{document}
